%; whizzy section -pdf -initex "pdflatex -ini"
\documentclass[a4paper,web]{frama-c-book}
\usepackage{hevea}
\usepackage{ifthen}

%%% Environnements dont le corps est suprim�, et
%%% commandes dont la d�finition est vide,
%%% lorsque PrintRemarks=false

\usepackage{comment}

\newcommand{\framac}{\textsc{Frama-C}\xspace}
\newcommand{\acsl}{\textsc{ACSL}\xspace}
\newcommand{\eacsl}{\textsc{E-ACSL}\xspace}
\newcommand{\C}{\textsc{C}\xspace}
\newcommand{\jml}{\textsc{JML}\xspace}

\newcommand{\nodiff}{\emph{No difference with \acsl.}}
\newcommand{\except}[1]{\emph{No difference with \acsl, but #1.}}
\newcommand{\limited}[1]{\emph{Limited to #1.}}
\newcommand{\absent}{\emph{No such feature in \eacsl.}}
\newcommand{\absentwhy}[1]{\emph{No such feature in \eacsl: #1.}}
\newcommand{\absentexperimental}{\emph{No such feature in \eacsl, since it is
    still experimental in \acsl.}}
\newcommand{\absentexcept}[1]{\emph{No such feature in \eacsl, but #1.}}
\newcommand{\difficultwhy}[2]{\emph{#1 is usually difficult to implement, since
    it requires #2. Thus you would not wonder if most tools do not support it
    (or support it partially).}}
\newcommand{\difficultswhy}[2]{\emph{#1 are usually difficult to implement,
    since they require #2. Thus you would not wonder if most tools do not
    support them (or support them partially).}}
\newcommand{\difficult}[1]{\emph{#1 is usually difficult to implement. Thus
    you would not wonder if most tools do not support it (or support
    it partially).}}
\newcommand{\difficults}[1]{\emph{#1 are usually difficult to implement. Thus
    you would not wonder if most tools do not support them (or support
    them partially).}}

\newcommand{\changeinsection}[2]{\textbf{Section \ref{sec:#1}:} #2.}

\newcommand{\todo}[1]{{\large \textbf{TODO: #1.}}}

\newcommand{\markdiff}[1]{{\color{blue}{#1}}}
\newenvironment{markdiffenv}[1][]{%
  \begin{changebar}%
  \markdiff\bgroup%
}%
{\egroup\end{changebar}}

% true = prints remarks for the ACSL working group.
% false = prints no remark for the distributed version of ASCL documents
\newboolean{PrintRemarks}
\setboolean{PrintRemarks}{false}

% true = prints marks signaling the state of the implementation
% false = prints only the ACSL definition, without remarks on implementation.
\newboolean{PrintImplementationRq}
\setboolean{PrintImplementationRq}{true}

% true = remarks about the current state of implementation in Frama-C
% are in color.
% false = they are rendered with an underline
\newboolean{ColorImplementationRq}
\setboolean{ColorImplementationRq}{true}

%% \ifthenelse{\boolean{PrintRemarks}}%
%%       {\newenvironment{todo}{%
%%             \begin{quote}%
%%             \begin{tabular}{||p{0.8\textwidth}}TODO~:\itshape}%
%%            {\end{tabular}\end{quote}}}%
%%       {\excludecomment{todo}}

\ifthenelse{\boolean{PrintRemarks}}%
      {\newenvironment{remark}[1]{%
             \begin{quote}\itshape%
             \begin{tabular}{||p{0.8\textwidth}}Remarque de {#1}~:}%
           {\end{tabular}\end{quote}}}%
      {\excludecomment{remark}}

\newcommand{\oldremark}[2]{%
\ifthenelse{\boolean{PrintRemarks}}{%
            %\begin{quote}\itshape%
            %\begin{tabular}{||p{0.8\textwidth}}Vieille remarque de {#1}~: #2%
            %\end{tabular}\end{quote}%
}%
{}}

\newcommand{\highlightnotreviewed}{%
\color{blue}%
}%

\newcommand{\notreviewed}[2][]{%
\ifthenelse{\boolean{PrintRemarks}}{%
  \begin{changebar}%
  {\highlightnotreviewed #2}%
  \ifthenelse{\equal{#1}{}}{}{\footnote{#1}}%
  \end{changebar}%
}%
{}}

\ifthenelse{\boolean{PrintRemarks}}{%
\newenvironment{notreviewedenv}[1][]{%
  \begin{changebar}%
  \highlightnotreviewed%
  \ifthenelse{\equal{#1}{}}{}{\def\myrq{#1}}%
  \bgroup}%
 {\egroup%
  \ifthenelse{\isundefined{\myrq}}{}{\footnote{\myrq}}\end{changebar}}}%
{\excludecomment{notreviewedenv}}

%%% Commandes et environnements pour la version relative � l'implementation
\newcommand{\highlightnotimplemented}{%
\ifthenelse{\boolean{ColorImplementationRq}}{\color{red}}%
           {\ul}%
}%

\newcommand{\notimplemented}[2][]{%
\ifthenelse{\boolean{PrintImplementationRq}}{%
  \begin{changebar}%
  {\highlightnotimplemented #2}%
  \ifthenelse{\equal{#1}{}}{}{\footnote{#1}}%
  \end{changebar}%
}%
{#2}}

\newenvironment{notimplementedenv}[1][]{%
\ifthenelse{\boolean{PrintImplementationRq}}{%
  \begin{changebar}%
  \highlightnotimplemented%
  \ifthenelse{\equal{#1}{}}{}{\def\myrq{#1}}%
  \bgroup
}{}}%
{\ifthenelse{\boolean{PrintImplementationRq}}{%
    \egroup%
    \ifthenelse{\isundefined{\myrq}}{}{\footnote{\myrq}}\end{changebar}}{}}

%%% Environnements et commandes non conditionnelles
\newcommand{\experimental}{\textsc{Experimental}}

\newsavebox{\fmbox}
\newenvironment{cadre}
     {\begin{lrbox}{\fmbox}\begin{minipage}{0.98\textwidth}}
     {\end{minipage}\end{lrbox}\fbox{\usebox{\fmbox}}}


\newcommand{\keyword}[1]{\lstinline|#1|\xspace}
\newcommand{\keywordbs}[1]{\lstinline|\\#1|\xspace}

\newcommand{\integer}{\keyword{integer}}
\newcommand{\real}{\keyword{real}}
\newcommand{\bool}{\keyword{boolean}}

\newcommand{\assert}{\keyword{assert}}
\newcommand{\terminates}{\keyword{terminates}}
\newcommand{\assume}{\keyword{assume}}
\newcommand{\requires}{\keyword{requires}}
\newcommand{\ensures}{\keyword{ensures}}
\newcommand{\exits}{\keyword{exits}}
\newcommand{\returns}{\keyword{returns}}
\newcommand{\breaks}{\keyword{breaks}}
\newcommand{\continues}{\keyword{continues}}
\newcommand{\assumes}{\keyword{assumes}}
\newcommand{\assigns}{\keyword{assigns}}
\newcommand{\reads}{\keyword{reads}}
\newcommand{\decreases}{\keyword{decreases}}

\newcommand{\boundseparated}{\keywordbs{bound\_separated}}
\newcommand{\Exists}{\keywordbs{exists}~}
\newcommand{\Forall}{\keywordbs{forall}~}
\newcommand{\bslambda}{\keywordbs{lambda}~}
\newcommand{\freed}{\keywordbs{freed}}
\newcommand{\fresh}{\keywordbs{fresh}}
\newcommand{\fullseparated}{\keywordbs{full\_separated}}
\newcommand{\distinct}{\keywordbs{distinct}}
\newcommand{\Max}{\keyword{max}}
\newcommand{\nothing}{\keywordbs{nothing}}
\newcommand{\numof}{\keyword{num\_of}}
\newcommand{\offsetmin}{\keywordbs{offset\_min}}
\newcommand{\offsetmax}{\keywordbs{offset\_max}}
\newcommand{\old}{\keywordbs{old}}
\newcommand{\at}{\keywordbs{at}}

\newcommand{\If}{\keyword{if}~}
\newcommand{\Then}{~\keyword{then}~}
\newcommand{\Else}{~\keyword{else}~}
\newcommand{\For}{\keyword{for}~}
\newcommand{\While}{~\keyword{while}~}
\newcommand{\Do}{~\keyword{do}~}
\newcommand{\Let}{\keywordbs{let}~}
\newcommand{\Break}{\keyword{break}}
\newcommand{\Return}{\keyword{return}}
\newcommand{\Continue}{\keyword{continue}}

\newcommand{\exit}{\keyword{exit}}
\newcommand{\main}{\keyword{main}}
\newcommand{\void}{\keyword{void}}

\newcommand{\struct}{\keyword{struct}}
\newcommand{\union}{\keywordbs{union}}
\newcommand{\inter}{\keywordbs{inter}}
\newcommand{\typedef}{\keyword{typedef}}
\newcommand{\result}{\keywordbs{result}}
\newcommand{\separated}{\keywordbs{separated}}
\newcommand{\sizeof}{\keyword{sizeof}}
\newcommand{\strlen}{\keywordbs{strlen}}
\newcommand{\Sum}{\keyword{sum}}
\newcommand{\valid}{\keywordbs{valid}}
\newcommand{\validrange}{\keywordbs{valid\_range}}
\newcommand{\offset}{\keywordbs{offset}}
\newcommand{\blocklength}{\keywordbs{block\_length}}
\newcommand{\baseaddr}{\keywordbs{base\_addr}}
\newcommand{\comparable}{\keywordbs{comparable}}

\newcommand{\N}{\ensuremath{\mathbb{N}}}
\newcommand{\ra}{\ensuremath{\rightarrow}}
\newcommand{\la}{\ensuremath{\leftarrow}}

% BNF grammar
\newcommand{\indextt}[1]{\index{#1@\protect\keyword{#1}}}
\newcommand{\indexttbs}[1]{\index{#1@\protect\keywordbs{#1}}}
\newif\ifspace
\newif\ifnewentry
\newcommand{\addspace}{\ifspace ~ \spacefalse \fi}
\newcommand{\term}[2]{\addspace\hbox{\lstinline|#1|%
\ifthenelse{\equal{#2}{}}{}{\indexttbase{#2}{#1}}}\spacetrue}
\newcommand{\nonterm}[2]{%
  \ifthenelse{\equal{#2}{}}%
             {\addspace\hbox{\textsl{#1}\ifnewentry\index{grammar entries!\textsl{#1}}\fi}\spacetrue}%
             {\addspace\hbox{\textsl{#1}\footnote{#2}\ifnewentry\index{grammar entries!\textsl{#1}}\fi}\spacetrue}}
\newcommand{\repetstar}{$^*$\spacetrue}
\newcommand{\repetplus}{$^+$\spacetrue}
\newcommand{\repetone}{$^?$\spacetrue}
\newcommand{\lparen}{\addspace(}
\newcommand{\rparen}{)}
\newcommand{\orelse}{\addspace$\mid$\spacetrue}
\newcommand{\sep}{ \\[2mm] \spacefalse\newentrytrue}
\newcommand{\newl}{ \\ & & \spacefalse}
\newcommand{\alt}{ \\ & $\mid$ & \spacefalse}
\newcommand{\is}{ & $::=$ & \newentryfalse}
\newenvironment{syntax}{\begin{tabular}{@{}rrll@{}}\spacefalse\newentrytrue}{\end{tabular}}
\newcommand{\synt}[1]{$\spacefalse#1$}
\newcommand{\emptystring}{$\epsilon$}
\newcommand{\below}{See\; below}

% colors

\definecolor{darkgreen}{rgb}{0, 0.5, 0}

% theorems

\newtheorem{example}{Example}[chapter]

\newcommand{\framacversion}{Nitrogen-20111001+dev
}

%Do not touch the following line. It is used in a Makefile hack to
%produce the ACSL documents for the ACSL working group.
%--\setboolean{PrintRemarks}{false}

%Do not touch the following line. It is used in a Makefile hack to
%produce the ACSL document shipped with the Frama-C distribution.
%--\setboolean{PrintImplementationRq}{false}

%\setboolean{ColorImplementationRq}{false}

%Already set in frama-c-book
%\usepackage[a4paper=true,pdftex,colorlinks=true,urlcolor=blue,pdfstartview=FitH]{hyperref}

%already set in frama-c-book
%\usepackage[T1]{fontenc}
%\usepackage{times}
\usepackage{amssymb}
\usepackage{graphicx}
%\usepackage{tikz}
\usepackage{color}
\usepackage{xspace}
\usepackage{makeidx}
\usepackage[normalem]{ulem}
\usepackage[leftbars]{changebar}
\usepackage{alltt}
\makeindex

%Already set in frama-c-book
%\setlength{\textheight}{240mm}
%\setlength{\topmargin}{-10mm}
%\setlength{\textwidth}{160mm}
%\setlength{\oddsidemargin}{0mm}
%\setlength{\evensidemargin}{0mm}
\newcommand{\version}{1.5+dev}

\renewcommand{\textfraction}{0.01}
\renewcommand{\topfraction}{0.99}
\renewcommand{\bottomfraction}{0.99}

% already taken care of in frama-c-book
% \usepackage{fancyhdr}
% \pagestyle{fancyplain}
% \renewcommand{\footrulewidth}{0.4pt}
% \addtolength{\headheight}{2pt}
% \addtolength{\headwidth}{1cm}
% \renewcommand{\chaptermark}[1]{\markboth{#1}{}}
% \renewcommand{\sectionmark}[1]{\markright{\thesection\ #1}}
% \lhead[\fancyplain{}{\bfseries\thepage}]{\fancyplain{}{\bfseries\rightmark}}
% \chead{}
% \rhead[\fancyplain{}{\bfseries\leftmark}]{\fancyplain{}{\bfseries\thepage}}
% \lfoot{\fancyplain{}{ANSI/ISO C Specification Language}}
% \cfoot{}
% \rfoot{\fancyplain{}{CAT RNTL project}}

\begin{document}
\sloppy
\hbadness=10000

\ifthenelse{\boolean{PrintImplementationRq}}%
  {\coverpage{\vbox{\mbox{ACSL Version \version}\newline
\mbox{Implementation in \framacversion{}}}}}%
  {\coverpage{\vbox{\mbox{ACSL: ANSI/ISO C Specification Language}\newline
        \mbox{Version \version}}}}

\begin{titlepage}
\includegraphics[height=14mm]{cealistlogo.jpg}
\hfill
\includegraphics[height=14mm]{inriasaclaylogo.png}
\vfill
\title{ACSL: ANSI/ISO C Specification Language}%
{Version \version{}\ifthenelse{\boolean{PrintImplementationRq}}%
{~--~\framacversion}{}}
\author{Patrick Baudin$^1$, Pascal Cuoq$^1$, Jean-Christophe Filli\^atre$^{4,3}$, Claude March\'e$^{3,4}$,\\ Benjamin Monate$^1$, Yannick Moy$^{2,4,3}$, Virgile Prevosto$^1$}

\begin{tabular}{l}
$^1$ CEA LIST, Software Reliability Laboratory, Saclay, F-91191 \\
$^2$ France T\'el\'ecom, Lannion, F-22307 \\
$^3$ INRIA Saclay - \^Ile-de-France, ProVal, Orsay, F-91893 \\
$^4$ LRI, Univ Paris-Sud, CNRS, Orsay, F-91405
\end{tabular}
\vfill
\begin{flushleft}
  \textcopyright 2009-2010 CEA LIST and INRIA

  This work has been supported by the `CAT' ANR project
  (ANR-05-RNTL-0030x) and by the ANR CIFRE contract 2005/973.
\end{flushleft}
\end{titlepage}

%%Contents should open on right
\cleardoublepage
\phantomsection
\label{chap:contents}
\tableofcontents

\chapter*{Foreword}

% This is a preliminary design of the ACSL language, a deliverable of
% the task 7.2 of the ANR RNTL project CAT
% (\url{http://www.rntl.org/projet/resume2005/cat.htm}). In this
% project, a reference implementation of ACSL is provided: the Frama-C
% platform (\url{http://frama-c.cea.fr}).

This is the version \version{} of ACSL design. Several features may still
evolve in the future. In particular, some features in this document
are considered \emph{experimental}, meaning that their syntax and
semantics is not already fixed.  These features are marked with
\experimental.  They must also be considered as advanced features,
which are not supposed to be useful for a basic use of that
specification language.

\section*{Acknowledgements}

We gratefully thank all the people who contributed to this document:
Sylvie Boldo,
Jean-Louis Cola\c{c}o,
Pierre Cr\'egut,
David Delmas,
St\'ephane Duprat,
Arnaud Gotlieb,
Thierry Hubert,
Dillon Pariente,
Pierre Rousseau,
Julien Signoles,
Jean Souyris.

%; whizzy-master "main.tex"
\chapter{Introduction}

This document is a reference manual for
\ifthenelse{\boolean{PrintImplementationRq}}%
{the \eacsl implementation provided by the Frama-C
  framework~\cite{frama-c}.}%
{E-ACSL.}
E-ACSL is an acronym for ``Executable ANSI/ISO C
Specification Language''. It is an ``executable'' subset of
\emph{stable} \acsl~\cite{acsl} implemented in the \framac
platform~\cite{framac}. ``Stable'' means that no experimental \acsl feature is
supported by \eacsl. Contrary to \acsl, each \eacsl specification is
executable: it may be evaluated at runtime.

In this document, we assume that the reader has a good knowledge of both
ACSL~\cite{acsl} and the ANSI C programming language~\cite{standardc99,KR88}.

\section{Organization of this document}

This document is organized in the very same way that the reference manual of
ACSL~\cite{acsl}.

Instead of being a fully new reference manual, this document points out the
differences between E-ACSL and ACSL. Each E-ACSL construct which is not pointed
out must be considered to have the very same semantics than its ACSL
counterpart. For clarity, each relevant grammar rules are given in BNF form
in separate figures like the \acsl reference manual does. In these rules,
constructs with semantic changes are displayed in \markdiff{blue}.

\section{Generalities about Annotations}\label{sec:gener-about-annot}
\nodiff

\section{Notations for grammars}
\nodiff


%; whizzy-master "main.tex"

\chapter{Specification language}
\label{chap:base}

\section{Lexical rules}

Lexical structure mostly follows that of ANSI C. A few differences
must be noted.
\begin{itemize}
\item The at sign (\verb|@|) is a blank character, thus equivalent
  to a space character.
\item Identifiers may start with the backslash character (\verb|\|).
\item Some UTF8 characters may be used in place of some constructs, as
  shown in the following table:

  \begin{center}
    \begin{tabular}{|l|l|l|}
      \hline
      \lstinline|>=| & \ensuremath{\geq} & 0x2265 \\
      \hline
      \lstinline|<=| & \ensuremath{\leq} & 0x2264 \\
      \hline
      \lstinline|>| & \ensuremath{>} & 0x003E \\
      \hline
      \lstinline|<| & \ensuremath{<} & 0x003C \\
      \hline
      \lstinline|!=| & \ensuremath{\not\equiv} & 0x2262 \\
      \hline
      \lstinline|==| & \ensuremath{\equiv} & 0x2261 \\
      \hline
      \lstinline|==>| & \ensuremath{\Longrightarrow} & 0x21D2 \\
      \hline
      \lstinline|<==>| & \ensuremath{\Longleftrightarrow} & 0x21D4 \\
      \hline
      \lstinline|&&| & \ensuremath{\land} & 0x2227 \\
      \hline
      \lstinline+||+ & \ensuremath{\lor} & 0x2228 \\
      \hline
      \lstinline+^^+ & \ensuremath{\oplus} & 0x22BB \\
      \hline
      \lstinline+!+ & \ensuremath{\neg} & 0x00AC \\
      \hline
      \lstinline+-+ (unary minus) & \ensuremath{-} & 0x2212 \\
      \hline
      \lstinline|\forall| & \ensuremath{\forall} & 0x2200 \\
      \hline
      \lstinline|\exists| & \ensuremath{\exists} & 0x2203 \\
      \hline
      \lstinline|integer| & \ensuremath{\mathbb{Z}} & 0x2124 \\
      \hline
      \lstinline|real| & \ensuremath{\mathbb{R}} & 0x211D \\
      \hline
      \lstinline|boolean| & \ensuremath{\mathbb{B}} & 0x1D539 \\
      \hline
    \end{tabular}
  \end{center}
\item Comments can be put inside ACSL annotations. They use the C++
  format, {\it i.e.} begin with \texttt{//} and extend to the end of
  current line.
\end{itemize}

%A bit complicated with listings
%In this document, we use UTF8 characters in the examples, to improve
%readability.

\section{Logic expressions}
\label{sec:expressions}

This first section presents the language of expressions one can use in
annotations. These are called \emph{logic expressions} in the following. They
correspond to pure C expressions, with additional constructs
that we will introduce progressively.

\begin{figure}[t]
  \begin{cadre}
    \input{term_modern.bnf}
  \end{cadre}
  \caption{Grammar of terms}
\label{fig:gram:term}
\end{figure}

\begin{figure}[hb]
  \begin{cadre}
    \input{predicate_modern.bnf}
  \end{cadre}
  \caption{Grammar of predicates}
\label{fig:gram:pred}
\end{figure}

Figures~\ref{fig:gram:term}~and~\ref{fig:gram:pred} present the
grammar for the basic constructs of logic expressions. In that
grammar, we distinguish between \emph{predicates}\index{predicate} and
\emph{terms}\index{term}, following the usual distinction between
propositions and terms in classical first-order logic. The grammar
for binders and type expressions is given separately in
Figure~\ref{fig:gram:binders}.

With respect to C pure expressions, the additional constructs are as follows:
\begin{description}
\item[Additional connectives] C operators \lstinline|&&| (UTF8:
  $\land$), \lstinline+||+ (UTF8: $\lor$) and \lstinline|!| (UTF8:
  $\neg$) are used as logical connectives. There are additional
  connectives \lstinline|==>| (UTF8: $\Longrightarrow$) for
  implication, \lstinline|<==>| (UTF8: $\Longleftrightarrow$) for
  equivalence and \lstinline|^^| (UTF8: $\oplus$) for exclusive
  or. These logical connectives all have a bitwise counterpart, either
  C ones like \lstinline|&|, \lstinline+|+, \lstinline|~| and
  \lstinline|^|, or additional ones like bitwise implication
  \lstinline|-->| and bitwise equivalence \lstinline|<-->|.

\item[Quantification] Universal quantification is denoted by
  \lstinline|\forall $\tau$ $x_1$,$\ldots$,$x_n$; e| and existential
  quantification by \lstinline|\exists $\tau$ $x_1$,$\ldots$,$x_n$; e|.

\begin{figure}[t]
  \fbox{\begin{minipage}{0.97\textwidth} \input{binders_modern.bnf}
    \end{minipage}}
  \caption{Grammar of binders and type expressions}
\label{fig:gram:binders}
\end{figure}

\item[Local binding]
    \lstinline|\let $x$ = $e_1$;$e_2$|
    introduces the name \lstinline|$x$| for
    expression \lstinline|$e_1$| which can be used in expression
    \lstinline|$e_2$|.

\item[Conditional] \lstinline|$c$ ? $e_1$ : $e_2$|. There is a subtlety
  here: the condition may be either a boolean term or a predicate.  In
  case of a predicate, the two branches must be also predicates, so
  that this construct acts as a connective with the following
  semantics: \lstinline|$c$ ? $e_1$ : $e_2$| is equivalent to
  \lstinline|($c$ ==> $e_1$) && (! $c$ ==> $e_2$)|.

\item[Syntactic naming] \lstinline|id : e| is a term or a predicate
  equivalent to $e$. It is different from local naming with \lstinline|\let|:
  the name cannot be reused in other terms or predicates. It is only
  for readibility purposes.

\item[Functional modifier]
 \begin{notimplementedenv}
  The composite element modifier is an additional
  operator related to C structure field and array accessors.
  The expression \lstinline|{ s \with .id = v }|
  denotes the same structure as \lstinline|s|, except for the field
  \lstinline|id| that is equal to \lstinline|v|.
  The equivalent expression for an array is
  \lstinline|{ t \with [ i ] = v }|
  which returns the same array as \lstinline|t|, except for the
  \lstinline|i$^{th}$| element whose value
  is \lstinline|v|.
  See section~\ref{sec:func-dep} for an example of use of these operators.
 \end{notimplementedenv}

\item[Logic functions] Applications in terms and in propositions are not
applications of C functions, but of logic functions or predicates; see
Section~\ref{sec:logicspec} for detail.

\item[Consecutive comparison operators]
  The construct
\lstinline|$t_1$ $relop_1$ $t_2$ $relop_2$ $t_3$ $\cdots$ $t_k$| with
  several consecutive comparison operators is a shortcut for
  \lstinline|($t_1$ $relop_1$ $t_2$) && ($t_2$ $relop_2$ $t_3$) && $\cdots$|.
  It is required that the \lstinline|$relop_i$| operators must be in
  the same ``direction'', \emph{i.e.} they must all belong either to
  $\{$\lstinline|<|, \lstinline|<=|, \lstinline|==|$\}$ or to
  $\{$\lstinline|>|,\lstinline|>=|,\lstinline|==|$\}$. Expressions such as
  \lstinline|x < y > z| or \lstinline|x != y != z| are not allowed.

  To enforce the same interpretation as in C expressions, one may need
  to add extra parentheses: \lstinline|a == b < c| is equivalent
  to \lstinline|a == b && b < c|, whereas
  \lstinline |a == (b < c)| is
  equivalent to \lstinline|\let x = b < c; a == x|.
  This situation raises some issues, as in the example below.

\oldremark{Virgile}{Dans le dernier cas, on pourrait vouloir avoir x, y
  et z deux � deux distincts, et eventuellement un pr�dicat n-aire
\distinct(x,y,z) comme dans SMT-lib}

\end{description}

There is a subtlety regarding comparison operators: they are
predicates when used in predicate position, and boolean functions when
used in term position.
\begin{example}
Let us consider the following example:
\begin{listing-nonumber}
  int f(int a, int b) { return a < b; }
\end{listing-nonumber}
\begin{itemize}
\item the obvious postcondition \lstinline[style=c]|\result == a < b|
  is not the right one because it is actually a shortcut for
  \lstinline[style=c]|\result == a && a < b|.
\item adding parentheses results in a correct post-condition
  \lstinline[style=c]|\result == (a < b)|. 
Note however that there is an implicit
  conversion (see Sec.~\ref{sec:typing})
  from the \lstinline[style=c]|int| (the type of \lstinline[style=c]|\result|) to
  \lstinline[style=c]|boolean| (the type of \lstinline[style=c]|(a<b)|)
\item an equivalent post-condition, which does not rely on implicit
  conversion, is \lstinline[style=c]|(\result != 0) == (a<b)|. Both pairs of
  parentheses are mandatory.
\item \lstinline[style=c]|\result == (integer)(a<b)| is also acceptable because it compares
  two integers. The cast towards \lstinline[style=c]|integer| enforces
  \lstinline[style=c]|a<b| to be understood as a boolean term. Notice that a cast
  towards \lstinline[style=c]|int| would also be acceptable.
\item \lstinline[style=c]|\result != 0 <==> a < b| is acceptable because it is an
  equivalence between two predicates.
\end{itemize}

\end{example}

\subsection{Operators precedence}

The precedence of C operators is conservatively extended with
additional operators, as shown in Figure~\ref{fig:precedence}. In this
table, operators are sorted from highest to lowest priority. Operators
of same priority are presented on the same line.


\begin{figure}[t]
  \begin{center}
    \begin{tabular}{|l|l|l|}
      \hline
      class 	& associativity & operators \\
      \hline
      selection & left & \lstinline|[$\cdots$]| \lstinline|->| \lstinline|.| \\
      unary 	& right & \lstinline|!| \lstinline|~| \lstinline|+|
      \lstinline|-| \lstinline|*| \lstinline|&| \lstinline|(cast)|
      \lstinline|sizeof| \\
      multiplicative & left & \lstinline|*| \lstinline|/|  \lstinline|%| \\
      additive & left & \lstinline|+| \lstinline|-| \\
      shift 	& left & \lstinline|<<| \lstinline|>>| \\
      comparison & left & \lstinline|<| \lstinline|<=| \lstinline|>| \lstinline|>=| \\
      comparison & left & \lstinline|==| \lstinline|!=| \\
      bitwise and & left & \lstinline|&| \\
      bitwise xor & left & \lstinline|^| \\
      bitwise or & left & \lstinline+|+ \\
      bitwise implies & left & \lstinline+-->+ \\
      bitwise equiv & left & \lstinline+<-->+ \\
      connective and     & left & \lstinline|&&| \\
      connective xor & left & \lstinline+^^+ \\
      connective or & left & \lstinline+||+ \\
      connective implies & right & \lstinline|==>| \\
      connective equiv & left & \lstinline|<==>| \\
      ternary connective & right & \lstinline|$\cdots$?$\cdots$:$\cdots$| \\
      binding & left & \Forall{} \Exists{} \Let{} \\
      naming & right & \lstinline|:| \\
      \hline
    \end{tabular}
  \end{center}
  \caption{Operator precedence}
\label{fig:precedence}
\end{figure}

There is a remaining ambiguity between the connective
\lstinline|$\cdots$?$\cdots$:$\cdots$| and the labelling operator
\lstinline|:|. Consider for instance the
expression \lstinline|x?y:z:t|. The precedence table does not indicate
whether this should be understood as \lstinline|x?(y:z):t| or
\lstinline|x?y:(z:t)|.
Such a case must be considered as a syntax error, and should be fixed
by explicitly adding parentheses.

\subsection{Semantics}
\label{sec:twovaluedlogic}

The semantics of logic expressions in ACSL is based on
mathematical first-order
logic~\cite{wikipedia-fo-logic}. In
particular, it is a 2-valued logic with only total functions. Consequently,
expressions are never ``undefined''.
This is an important design choice and the specification writer should
be aware of that. (For a discussion about the issues raised by such
design choices, in similar specification languages such as JML, see
the comprehensive list compiled by Patrice
Chalin~\cite{chalin05ftfjp,chalin07icse}.)

Having only total functions implies than one can write
terms such as \lstinline|1/0|, or \lstinline|*p| when \lstinline!p! is
null (or more
generally
when it points to a non-properly allocated memory cell). In
particular, the predicates
\begin{tabular}{rcl}
  \lstinline!1/0! &\lstinline|==|& \lstinline!1/0! \\\relax
  \lstinline!*p! &\lstinline|==|& \lstinline!*p!
\end{tabular}
are valid, since they are instances of the axiom $\forall x,
x=x$ of first-order logic. The reader should not be alarmed, because
there is no way to deduce anything useful from such terms.
\oldremark{Anonyme}{Donc on ne peut pas introduire de construction
non-d�terministe comme  $(\texttt{any} x \mid P)$
}
As usual, it is up to the specification designer to write
consistent assertions. For example, when introducing the following lemma
(see Section~\ref{sec:logicspec}):
\listinginput{1}{div_lemma.c}
a premise is added to require \lstinline|y| to be non zero.


\subsection{Typing}\label{sec:typing}

The language of logic expressions is typed (as in \emph{multi-sorted}
first-order logic). Types are either C types or \emph{logic types}
defined as follows:
\begin{itemize}
\item ``mathematical'' types: \lstinline|integer|\indextt{integer} for
  unbounded, mathematical integers, \lstinline|real|\indextt{real} for
  real numbers, \lstinline|boolean|\indextt{boolean} for booleans (with
  values written \lstinline|\true|\indexttbs{true} and
  \lstinline|\false|\indexttbs{false});
\item logic types introduced by the specification writer (see
  Section~\ref{sec:logicspec}).
\end{itemize}
There are implicit coercions for numeric types:
\begin{itemize}
\item C integral types \lstinline|char|, \lstinline|short|, \lstinline|int| and
  \lstinline|long|, signed or unsigned, are all subtypes of type
\lstinline|integer|;
\item \lstinline|integer| is itself a subtype of type \lstinline|real|;
\item C types \lstinline|float| and \lstinline|double| are subtypes of
  type \lstinline|real|.
\end{itemize}
Notes:
\begin{itemize}
\item There is a distinction between booleans and predicates. The
  expression \lstinline|x<y| in term position is a boolean, and the
  same expression
  is also allowed in predicate position.
\item Unlike in C, there is a distinction between booleans and
  integers. There is an implicit promotion from integers to booleans,
  thus one may write \lstinline|x && y| instead of \lstinline|x != 0 && y != 0|.
  If the reverse conversion is needed, an explicit cast\index{cast} is
  required, \emph{e.g.} \lstinline|(int)(x>0)+1|, where
  \lstinline|\false| becomes 0
  and \lstinline|\true| becomes 1.
\item Quantification can be made over any type: logic types and C
  types. Quantification over pointers must be used carefully,
  since it depends on the memory state where dereferencing is done (see
  Section~\ref{sec:quantification} and
  Section~\ref{sec:logicalstates}).
\end{itemize}

Formal typing rules for terms are given in appendix~\ref{sec:typingrules}.

\subsection{Integer arithmetic and machine integers}

The following integer arithmetic operations apply to
\emph{mathematical integers}: addition, subtraction, multiplication,
unary minus. The value of a C variable of an integral type is promoted
to a mathematical integer. As a consequence, there is no
``arithmetic overflow'' in logic expressions.

Division and modulo are also mathematical operations, which coincide
with the corresponding C operations on C machine integers, thus
following the ISO C99 conventions. In particular, these are not the
usual mathematical Euclidean division and remainder.
Generally speaking, division
rounds the result towards zero. The results are not specified if
divisor is zero; otherwise if $q$ and $r$ are the quotient and the
remainder of $n$ divided by $d$ then:
\begin{itemize}
\item $|d\times q| \leq |n|$, and $|q|$ is maximal for this property;
\item $q$ is zero if $|n|<|d|$;
\item $q$ is positive if $|n|\geq|d|$ and $n$ and $d$ have the same sign;
\item $q$ is negative if $|n|\geq|d|$ and $n$ and $d$ have opposite signs;
\item $q\times d+r = n$;
\item $|r|<|d|$;
\item $r$ is zero or has the same sign as $n$.
\end{itemize}

\begin{example}
  The following examples illustrate the results of division and modulo
  depending on the sign of their arguments:
  \begin{itemize}
  \item \lstinline|5/3| is 1 and \lstinline|5%3| is 2;
  \item \lstinline|(-5)/3| is -1 and \lstinline|(-5)%3| is -2;
  \item \lstinline|5/(-3)| is -1 and \lstinline|5%(-3)| is 2;
  \item \lstinline|(-5)/(-3)| is 1 and \lstinline|(-5)%(-3)| is -2.
  \end{itemize}
\end{example}

\subsubsection{Hexadecimal and octal constants}

Hexadecimal and octal constants are always non-negative. Suffixes
\texttt{u} and \texttt{l} for C constants are allowed but meaningless.

\subsubsection{Casts and  overflows}\index{cast}

In logic expressions, casting from mathematical integers
to an integral C type \lstinline|t| (such as \lstinline|char|,
\lstinline|short|, \lstinline|int|, etc.) is allowed and is
interpreted as follows: the result is the unique value of the
corresponding type that is congruent to the mathematical result modulo
the cardinal of this type, that is $2^{8\times\sizeof(t)}$.
\begin{example}
  \lstinline|(unsigned char)1000| is $1000 \bmod 256$ i.e. $232$.
\end{example}

To express in the logic the value of a C
expression, one has to add all the necessary casts. For
example, the logic expression denoting the value of the C expression
\lstinline|x*y+z| is \lstinline|(int)((int)(x*y)+z)|.
\oldremark{Claude}{Si la conjecture suivante est vraie, alors il
  faudrait en parler: Si e1 et e2 sont des expressions d'un type
  entier de C, quels qu'ils soient et si op est une operation binaire
  sur entiers, alors l'operation C $e1~op~e2$ (de type $\tau$),donne
  le meme resultat que l'expression logique $(\tau)(e_1~op~e_2)$ ou
  $op$ est l'operation logique sur $Z$. Ce cast permet donc a
  l'utilisateur de parler du resultat du calcul C d'une expressions
  dans la logique
}
Note that there is no implicit cast from integers to C integral types.
\begin{example}
  The declaration
  \begin{listing-nonumber}
    //@ logic int f(int x) = x+1 ;
  \end{listing-nonumber}
  is not allowed because \lstinline!x+1!, which is a mathematical
  integer, must be cast to \lstinline|int|.  One should write either
  \begin{listing-nonumber}
    //@ logic integer f(int x) = x+1 ;
  \end{listing-nonumber}
  or
  \begin{listing-nonumber}
    //@ logic int f(int x) = (int)(x+1) ;
  \end{listing-nonumber}
\end{example}

\subsubsection{Quantification on C integral types}
\label{sec:quantification}

Quantification over a C integral type corresponds to integer
quantification over the corresponding interval.
\begin{example}
Thus the formula
\begin{listing-nonumber}
\forall char c; c <= 1000
\end{listing-nonumber}
is equivalent to
\begin{listing-nonumber}
\forall integer c; CHAR_MIN <= c <= CHAR_MAX ==> c <= 1000
\end{listing-nonumber}
where the the bounds \verb|CHAR_MIN| and \verb|CHAR_MAX| are defined
in \verb|limits.h|
\end{example}


\subsubsection{Size of C integer types}

The size of C types is architecture-dependent. ACSL does
not enforce these sizes either, hence the semantics of terms involving
such types is also architecture-dependent. The
\lstinline|sizeof|\indextt{sizeof} operator may be used in annotations and is
consistent with its C counterpart. For instance, it should be
possible to verify the following code:

\listinginput{1}{sizeof.c}

Constants giving maximum and minimum values of those types may be
provided in a library.


\subsubsection{Enum types}

Enum types are also interpreted as mathematical integers. Casting\index{cast}
an integer into an enum in the logic gives the same result as
if the cast was performed in the C code.

\subsubsection{Bitwise operations}

Like arithmetic operations, bitwise operations apply to any
mathematical integer: any mathematical integer has a unique infinite
2-complement binary representation with infinitely many zeros (for
non-negative numbers) or ones (for negative numbers) on the left.
Bitwise operations apply to this representation.
\begin{example} ~
  \begin{itemize}
  \item \lstinline|7 & 12 == $\cdots$00111 & $\cdots$001100 == $\cdots$00100 == 4|
  \item \lstinline+-8 | 5 == $\cdots$11000 | $\cdots$00101 == $\cdots$11101 == -3+
  \item \lstinline+~5 == ~$\cdots 00101$ == $\cdots$111010 == -6+
  \item \lstinline+-5 << 2 == $\cdots$11011 << 2 == $\cdots$11101100 == -20+
  \item \lstinline+5 >> 2 == $\cdots$00101 >> 2 == $\cdots$0001 == 1+
  \item \lstinline+-5 >> 2 == $\cdots$11011 >> 2 == $\cdots$1110 == -2+
  \end{itemize}
\end{example}

\subsection{Real numbers and floating point numbers}

Floating-point constants and operations are interpreted as
mathematical real numbers: a C variable of type float or double is
implicitly promoted to a real. Integers are promoted to reals if
necessary. Usual binary operations are interpreted as operators on
real numbers, hence they never involve any rounding nor overflow.

\begin{example}
  In an annotation, \lstinline|1e+300 * 1e+300| is equal to
  \lstinline|1e+600|, even if that last number exceeds the largest
  representable number in double precision: there is no "overflow".

  $2 * 0.1$ is equal to the real number $0.2$, and not to any of its
  floating-point approximation: there is no "rounding".
\end{example}
Unlike the promotion of C integer types to mathematical integers,
there are special float values which do not naturally map to a real
number, namely the IEEE-754 special values for ``not-a-number'',
$+\infty$ and $-\infty$. See below for a detailed discussion on such
special values. However, remember that ACSL's logic has only total
functions. Thus, there are implicit promotion functions
\lstinline|real_of_float|\indextt{real\_of\_float} and
\lstinline|real_of_double|\indextt{real\_of\_double} whose results on
the 3 values above is left unspecified.

In logic, real literals can also be expressed under the hexadecimal
form of C99: \texttt{0x$hh$.$hh$p$\pm{}dd$} where $h$ are hexadecimal
digits and $dd$ is in decimal, denotes number $hh.hh\times
2^{dd}$, e.g. \texttt{0x1.Fp-4} is $(1+15/16)\times 2^{-4}$.


Usual operators for comparison are interpreted as real operators
too. In particular, equality operation $\equiv$ of float (or double)
expressions means equality of the real numbers they represent
respectively. Or equivalently, $x \equiv y$ for $x,y$ two float
variables means \lstinline|real_of_float(x) $\equiv$ real_of_float(y)|
with the mathematical equality of real
numbers.


Special predicates are also available to express the comparison
operators of float (resp. double) numbers as in C:
\lstinline|\eq_float|\indexttbs{eq\_float}\indexttbs{eq\_double},
\lstinline|\gt_float|\indexttbs{gt\_float}\indexttbs{gt\_double},
\lstinline|\ge_float|\indexttbs{ge\_float}\indexttbs{ge\_double},
\lstinline|\le_float|\indexttbs{le\_float}\indexttbs{le\_double},
\lstinline|\lt_float|\indexttbs{lt\_float}\indexttbs{lt\_double},
\lstinline|\ne_float|\indexttbs{ne\_float}\indexttbs{ne\_double}
(resp. for double).



\subsubsection{Casts, infinity and NaNs}

Casting\index{cast} from a C integer type or a float type to a float
or a double is as in C: the same conversion operations apply.

Conversion of real numbers to float or double values indeed depends on
various possible rounding modes defined by the IEEE~754
standard~\cite{float-std,wikipedia-float}. These modes are defined by
a logic type (see section~\ref{sec:concrete-logic-types}):
\begin{listing-nonumber}
/*@ type rounding_mode = \Up | \Down | \ToZero | \NearestAway | \NearestEven;
*/
\end{listing-nonumber}
Then rounding a real number can be done explicitly using functions
\begin{listing-nonumber}
  logic float \round_float`\indexttbs{round\_float}`(rounding_mode m, real x);
  logic double \round_double`\indexttbs{round\_double}`(rounding_mode m, real x);
\end{listing-nonumber}
Cast operators \lstinline|(float)| and \lstinline|(double)| applied to a
mathematical integer or real number $x$ are equivalent to apply
rounding functions above with the nearest-even rounding mode (which is
the default rounding mode in C programs)
%(see IEEE 754 norm), which is the closest to $x$ real number
%which is representable in the float (resp. double) type.
If the source real number is too large, this may also result into one
of the special values +infinity and -infinity.
\begin{example}
  We have
  \lstinline|(float)0.1| $\equiv 13421773 \times 2 ^{-27}$ which is equal to
  $0.100000001490116119384765625$
\end{example}
Notice also that unlike for integers, suffixes \lstinline|f| and
\lstinline|l| are
meaningful, because they implicitly add a cast operator as above.

% \begin{remark}{Claude}
%   ne devrait-on pas ajouter le suffixe d dans la logique ?

%   par ex. 0.1d -> (double)0.1

% \end{remark}

This semantics of casts ensures that the float result \lstinline|r| of
a C operation \lstinline|$e_1$ op $e_2$| on floats, if there is no
overflow and if the default
rounding mode is not changed in the program, has the same real value
as the logic expression \lstinline|(float)($e_1$ op $e_2$)|. Notice
that this is not
true for the equality \lstinline|\eq_float| of floats: \lstinline|-0.0 + -0.0|
in C is equal to the float number \lstinline|-0.0|, which is not
\lstinline|\eq_float| to \lstinline|0.0|, which is the value of the logic
expression \lstinline|(float)(-0.0 + -0.0)|.

Finally, additional predicates are provided on float and double
numbers, which check that their argument is NaN or a
finite number respectively:
\begin{listing}{1}
predicate \is_NaN`\indexttbs{is\_NaN}`(float x);
predicate \is_NaN`\indexttbs{is\_NaN}`(double x);
predicate \is_finite`\indexttbs{is\_finite}`(float x);
predicate \is_finite`\indexttbs{is\_finite}`(double x);
\end{listing}

\begin{remark}{Claude}

  et les long double

\end{remark}

\oldremark{Claude}{

on pourrait enonce un certain nombre de proprietes:

\is_finite_float(x) ==> (float)(\real_of_float(x)) == x

}

\subsubsection{Quantification}

Quantification over a variable of type \lstinline|real| is of course usual
quantification over real numbers.

Quantification over float (resp. double) types is allowed too, and is
supposed to range over all real numbers representable as floats (resp
doubles). In particular, this does not include NaN, +infinity and
-infinity in the considered range.


\subsubsection{Mathematical functions}

Classical mathematical operations like exponential, sine, cosine, and
such are available as built-in:

\begin{remark}{Virgile}
  Il faut indexer ces fonctions

  Il faut mentionner quelque part que la surcharge des finctions
  logiques est possible.

\end{remark}

\begin{listing-nonumber}
integer \min(integer x, integer y) ;
integer \max(integer x, integer y) ;
real \min(real x, real y) ;
real \max(real x, real y) ;
integer \abs(integer x) ;
real \abs(real x) ;

real \sqrt(real x) ;
real \pow(real x, real y) ;

integer \ceil(real x) ;
integer \floor(real x) ;

real \exp(real x) ;
real \log(real x) ;
real \log10(real x) ;

real \cos(real x) ;
real \sin(real x) ;
real \tan(real x) ;

real \cosh(real x) ;
real \sinh(real x) ;
real \tanh(real x) ;

real \acos(real x) ;
real \asin(real x) ;
real \atan(real x) ;

real \atan2(real y, real x) ;
real \hypot(real x, real y) ;
\end{listing-nonumber}


\oldremark{All}{  syntaxe pour abs et power ? a priori non, on peut
  fournir une library pour les reels avec abs(), power(), mais
  proposer des syntaxes speciales parait difficile

  mais: on pourrait permettre des definition d'operateurs infixes

}

\subsubsection{Exact computations}

In order to specify properties of rounding errors, it is useful to express
something about the so-called \emph{exact}
computations~\cite{BoldoFilliatre07}: the computations that would be performed
in an ideal mode where variables denote true real numbers.

To express such exact computations, two special
constructs exist in annotations:
\begin{itemize}
\item $\verb|\exact|(x)$ denotes the value of the C variable $x$ (or
  more generally any C left-value) as if the program was executed with
  ideal real numbers.
\item $\verb|\round_error|(x)$ is a shortcut for $|x-\verb|\exact|(x)|$
\end{itemize}

\begin{example}
Here is an example of a naive approximation of cosine~\cite{ayad09}.
\begin{listing-nonumber}
/*@ requires \abs(\exact(x)) <= 0x1p-5;
  @ requires \round_error(x) <= 0x1p-20;
  @ ensures \abs(\exact(\result) - \cos(\exact(x))) <= 0x1p-24;
  @ ensures \round_error(\result) <= \round_error(x) + 0x3p-24;
  @*/
float cosine(float x) {
  return 1.0f - x * x * 0.5f;
}
\end{listing-nonumber}
\end{example}

\begin{remark}{JCF+All}

Que deviennent les champs modeles de caduceus ?

Il faudrait les integrer a une traitement plus general des modeles
pour des types de donnees

\end{remark}




\subsection{C arrays and pointers}

\subsubsection{Address operator, array access, pointer arithmetic and dereferencing}
\label{sec:address}

These operators are similar to their corresponding C operators.

\emph{address-of} operator should be used with caution.
Values in logic do not lie in C memory so it does not mean
anything to talk about their ``address''.

\begin{remark}{???}
The type system should provide a boolean information on whether a
given logic term denotes a C memory location (i.e. a C left-value).
\end{remark}

Unlike in C, 
there is no implicit cast from from an array type to a pointer type.
Nevertheless, arithmetic and dereferencing over arrays lying in C memory are 
allowed like in C.

\begin{example}
Dereferencing a C array is equivalent to an access 
to the first element of the array~;
shifting it from $i$ denotes the address of its $i^{th}$ element.
\begin{listing-nonumber}
int tab[10] = { 1 } ;
int x ;
int *p = &x;

//@ requires p == &x
int main(void){
  //@ assert tab[0]==1 && *p == x;
  //@ assert *tab == 1;
  int *q = &tab[3];
  //@ assert q+1 == tab+4;
  ...
}
\end{listing-nonumber}
\end{example}

Since pointers can only refers to values lying in C memory, 
\lstinline|p->s| is always equivalent to \lstinline|(*p).s|.
On the contrary, \lstinline|t[i]| is not always equivalent to 
\lstinline|*(t+i)|, especially for arrays not lying in C memory. 
Section~\ref{sec:aggregate} details the use of arrays as logic values. There
are also differences between \lstinline|t| and the pointer to its first element
when evaluating an expression at a given program point. See 
Section~\ref{sec:at} for more information.

\subsubsection{Function pointers}

Pointers to C functions are allowed in logic. The only possible use of
them is to check for equality.

\begin{example}
~
\begin{listing-nonumber}
int f(int x);
int g(int x);

//@ requires p == &f || p == &g;
void h(int(*p)(int)) {
...
}
\end{listing-nonumber}
\end{example}

\subsection{Structures, Unions and Arrays in logic}
\label{sec:aggregate}

Aggregate C objects (i.e. structures, unions and arrays) are also
possible values for terms in logic. They can be passed as parameters (and
also returned) to logic functions, tested for equality, etc. like any
other values.

Aggregate types can be declared in logic, 
\notimplemented{and their contents may be
any logic types themselves}. Constructing such values in logic can be
performed using a syntax similar to C designated initializers.

\begin{example}
Array types in logic may be declared either with or without an
explicit non-negative length. Access to the length of a logic array
can be done with \lstinline|\length|.

\begin{listing-nonumber}
//@ type point = struct { real x; real y; };
//@ type triangle = point[3];

//@ logic point origin = { .x = 0.0 , .y = 0.0 };
/*@ logic triangle t_iso = { [0] = origin,
  @                          [1] = { .y = 2.0 , .x = 0.0 }
  @                          [2] = { .x = 2.0 , .y = 0.0 }};
  @*/

/*@ logic point centroid(triangle t) = {
  @    .x = mean3(t[0].x,t[1].x,t[2].x);
  @    .y = mean3(t[0].y,t[1].y,t[2].y);
  @ };
  @*/

//@ type polygon = point[];
/*@ logic perimeter(polygon p) =
  @   \sum(0,\length(p)-1,\lambda integer i;d(p[i],p[(i+1) % \length(p)])) ;
  @*/
\end{listing-nonumber}
Beware that because of the principle of only total functions in logic,
t[i] can appear in ACSL annotations even if i is outside the array bounds.

\end{example}

\subsubsection{Functional updates}

Syntax for functional update is similar to initialization of aggregate objects.
\begin{example}
Functional update of an array is done by
\begin{listing-nonumber}
{ t_iso \with [0] = { .x = 3.0, .y = 3.0 } }
\end{listing-nonumber}

Functional update of a structure is done by
\begin{listing-nonumber}
{ origin \with .x = 3.0 }
\end{listing-nonumber}

There is no particular syntax for functional update of an union. 
For an \lstinline|object| of an union type, the following equality is not true 
\begin{listing-nonumber}
{ { object \with .x = 3.0 }
           \with .y = 2.0 } == { { object \with .y = 2.0 }
                                          \with .x = 3.0 }
\end{listing-nonumber}

\end{example}

\begin{remark}{}
other built-in constructs: array.sub ? array.map ? (later)
\end{remark}



The equality predicate \lstinline|==| applies to aggregate values, but it
is required that they have the same type. Then equality amounts the
recursively check equality of fields.  Equality of arrays of different
lengths returns false.  Beware that equality of unions is also
equality of all fields.


\subsubsection{C aggregate types}

C aggregates types (struct, union or array) naturally map to logic
types, by recursively mapping their fields.

\begin{example}
There is no implicit cast to type of the updated/initialized fields.
\begin{listing-nonumber}
struct S { int x; float y; int t[10]; };

//@ logic integer f(struct S s) = s.t[3];
//@ logic struct S g(integer n, struct S s) = { s \with .x = (int)n };
\end{listing-nonumber}
Unlike in C, all fields should be initialized:
\begin{listing-nonumber}
/*@ logic struct S h(integer n, int a[10]) = {
  @   .x = (int)n, .y = (float)0.0, .t = a
  @   };
  @*/
\end{listing-nonumber}
\end{example}

\subsubsection{Cast and conversion}

Unlike in C, there is no implicit conversion from an array type to a
pointer type. On the other hand, there is an implicit conversion from
an array of a given size to an array with unspecified size (but not
the converse).

\begin{example}
~
\begin{listing-nonumber}
//@ logic point square[4] = { origin, ... };

//@ ... perimeter(square);          // well typed
//@ ... centroid(square);           // wrongly typed
//@ ... centroid((triangle)square); // well-typed (truncation)
\end{listing-nonumber}
\end{example}

An explicit cast from an array type to a pointer type is allowed only
for arrays that lies in C memory. As in C, the result of the cast is
the address of the first element of the array (see
Section~\ref{sec:address}).

Conversely, an explicit cast from a pointer type to an array type or a
structure type is allowed, and acts as collecting the values it points
to.

Subtyping and cast recursively applies to fields.

\begin{example}
~
\begin{listing-nonumber}
  struct { float u,v; } p[10];

  //@ assert centroid((point[3])p) == ...

  //@ assert perimeter((point[])p) == ...
\end{listing-nonumber}
\end{example}

Precisely, conversion of a pointer $p$ of type $\tau*$ to an logic array of type
$\tau[]$ returns a logic array $t$ such that
\[
\lstinline|length|(t) = (
\blocklength\indexttbs{block\_length}
(p) - 
\lstinline|\offset|
(p))/\lstinline|sizeof|(\tau)
\]

\subsection{String literals}

\begin{remark}{Pascal, Benjamin}

Quelle est la semantique des constanes chaines dans les annotations
logiques ?

Une reponse simple possible : equivalent a une constante globale ghost


"foo" devient:


ghost char string\_literal\_foo [4] = "foo" (avec le 0 final)

et utiliser string\_literal\_foo a la place de "foo"


\end{remark}



\section{Function contracts}
\label{sec:fn-behavior}
\index{function contract}\index{contract}

\begin{figure}[t]
  \begin{cadre}
      \input{fn_behavior_modern.bnf}
   \end{cadre}
    \caption{Grammar of function contracts}
  \label{fig:gram:contracts}
\end{figure}

\begin{figure}[t]
  \begin{cadre}
      \input{oldandresult_modern.bnf}
    \end{cadre}
    \caption{\protect\old and \protect\result in terms}
  \label{fig:gram:oldandresult}
\end{figure}

Figure~\ref{fig:gram:contracts} shows a grammar for function
contracts. \textsl{Location} denotes a memory location and is defined
in Section~\ref{sec:locations}. 
% [Patrick B.] START proposition:
\textsl{Allocation-clauses} allow to specify which memory location 
is dynamically allocated or deallocated by the function from the \textsl{heap}; 
they are defined later in Section~\ref{sec:allocation-clause}.
% END proposition

This section is organized as follows.  First, the grammar for terms is
extended with two new constructs.  Then
Section~\ref{sec:simplecontracts} introduces \emph{simple contracts}.
Finally, Section~\ref{subsec:behaviors} defines more general contracts
involving \emph{named behaviors}. 

The \lstinline|decreases| and
\lstinline|terminates| clauses are presented later in
Section~\ref{sec:termination}.
\textsl{Abrupt-clauses}
allow to specify what happens when the function does not return
normally but \lstinline|exit|s abruptly; they are defined in
Section~\ref{sec:abrupt-clause}.


\subsection{Built-in constructs %
  \texorpdfstring{\old}{\textbackslash{}old} %
 and \texorpdfstring{\result}{\textbackslash{}result}}

Post-conditions usually require to refer to both the function result and
values in the pre-state. Thus terms are extended with the following new
constructs (shown in Figure~\ref{fig:gram:oldandresult}).
\begin{itemize}
\item \lstinline|\old(e)|\indexttbs{old} denotes the value of
  predicate or term \lstinline|e| in the pre-state of the function.
\item \result\indexttbs{result} denotes the returned value of the function.
\end{itemize}
\lstinline|\old(e)| and \result{} can be used only in \ensures{} clauses, since
the other clauses already refer to the pre-state. In addition,
\result{} can not be used in the contract of a function which returns
\void{}.
% See Section~\ref{sec:at} for more details on scoping rules for labels.

C function parameters\index{formal parameter} are obtained by
value from actual parameters that mostly remain unaltered by the
function calls. For that reason, formal parameters in function
contracts are defined such that they always refer implicitly to their
values interpreted in the pre-state.
Thus, \old construct is useless for formal parameters (in
function contracts only).

\subsection{Simple function contracts}
\label{sec:simplecontracts}

A simple function contract, having only simple clauses and no named behavior, takes the following
form:\indextt{requires}\indextt{assigns}\indextt{ensures}
\begin{listing}{1}
/*@ requires P$_1$; requires P$_2$; `\dots`
  @ assigns L$_1$;  assigns L$_2$;  `\dots`
  @ ensures E$_1$;  ensures E$_2$;  `\dots`
  @*/
\end{listing}
The semantics of such a contract is as follows:
\begin{itemize}
\item The caller of the function must guarantee that it is called in a
  state where the property \lstinline|P$_1$ && P$_2$ && ...| holds.
\item The called function returns a state where the property
  \lstinline|E$_1$ && E$_2$ && ...| holds.
\item All memory locations that are allocated in both the pre-state and 
  the post-state\footnote{Functions that allocate or free memory can be 
    specified with additional clauses described in
    section~\ref{sec:allocation-clause}.} 
  and do not belong to the
  set \lstinline|L$_1$ $\cup$ L$_2$ $\cup \dots$| are left unchanged in the
  post-state. The set \lstinline|L$_1$ $\cup$ L$_2$ $\cup \dots$| 
  itself is interpreted in the pre-state.
\end{itemize}

% TODO : donner une vision plus sym�trique appelant/appel� d'un contrat

Having multiple \lstinline|requires|, \lstinline|assigns|, or 
\lstinline|ensures| clauses only improves 
readibility since the contract above is equivalent to the following
simplified one:
\begin{listing}{1}
/*@ requires P$_1$ && P$_2$ && `\dots`;
  @ assigns L$_1$, L$_2$,`\dots`;
  @ ensures E$_1$ && E$_2$ && `\dots`;
  @*/
\end{listing}
If no clause \requires\ is given, it defaults to
\lstinline!\true!, and similarly for \ensures\ clause.
Giving no \assigns\ clause
means that locations assigned by the function are not specified, so
the caller has no information at all on this function's side
effects. See Section~\ref{sec:multiplecontracts} for more details on
default status of clauses.

\oldremark{Colaco}{

  les assigns doivent etre exhaustifs, et doivent etre prouve corrects
  selon le corps de la function.

}

\begin{example}
  The following function is given a simple contract for computation of
  the integer square root.

  \listinginput{1}{isqrt.c}
%
  The contract means that the function must be called with a
  nonnegative argument, and returns a value satisfying
  the conjunction of the three \ensures{} clauses.
  Inside these \ensures{} clauses, the use of the construct \lstinline|\old(x)|
  is not necessary, even if the function modifies the formal
  parameter\index{formal parameter}
  \lstinline|x|, because function calls modify a copy of the effective parameters,
  and the effective parameters remain unaltered.  In fact, \lstinline|x| denotes
  the effective parameter of \lstinline|isqrt| calls which has the same value
  interpreted in the pre-state than in the post-state.

\end{example}

\begin{example}
  The following function is given a contract to specify that it increments
  the value pointed to by the pointer given as argument.

  \listinginput{1}{incrstar.c}
%
  The contract means that the function must be called with a pointer
  \lstinline|p| that points to a safely allocated memory location (see
  Section~\ref{sec:pointers} for details on the \valid built-in
  predicate). It does not modify any memory location but the one
  pointed to by \lstinline|p|. Finally, the \ensures clause specifies that
  the value \lstinline!*p! is incremented by one.
\end{example}

\subsection{Contracts with named behaviors}
\label{subsec:behaviors}
\index{function behavior}\index{behavior}
The general form of a function contract contains several named
behaviors (restricted to two behaviors, in the following, for
readability).
\lstset{firstnumber=auto}
\begin{lstlisting}[style=c,basicstyle=\lp@basic,numbers=left,name=behaviors]
/*@ requires P;
  @ behavior b$_1$:
  @   assumes  A$_1$;
  @   requires R$_1$;
  @   assigns  L$_1$;
  @   ensures  E$_1$;
  @ behavior b$_2$:
  @   assumes  A$_2$;
  @   requires R$_2$;
  @   assigns  L$_2$;
  @   ensures  E$_2$;
  @*/
\end{lstlisting}
The semantics of such a contract is as follows:
\begin{itemize}
\item The caller of the function must guarantee that the call is
  performed in a state where the property
  \lstinline|P && (A$_1$ ==> R$_1$) && (A$_2$ ==> R$_2$)| holds.
\item The called function returns a state where
the properties \lstinline|\old(A$_i$) ==> E$_i$| hold for each~$i$.
\item For each $i$, if the function is called in a pre-state where
  \lstinline|A$_i$| holds, then each memory location of that pre-state 
  that does not belong to the set \lstinline|L$_i$| is left unchanged in 
  the post-state.
% In any case, locations of $L$ are unchanged.
\end{itemize}

\requires{} clauses in the behaviors are proposed
mainly to improve readibility (to avoid some duplication of
formulas), since the contract above is equivalent to the following
simplified one:
\begin{listing}{1}
/*@ requires P && (A$_1$ ==> R$_1$) && (A$_2$ ==> R$_2$);
  @ assigns L$_1$,L$_2$;
  @ ensures \old(A$_1$) ==> E$_1$ && \old(A$_2$) ==> E$_2$;
  @*/
\end{listing}

A simple contract such as
\begin{listing}{1}
/*@ requires P; assigns L; ensures E; */
\end{listing}
is actually equivalent to a single named behavior as follows:
\begin{listing}{1}
/*@ requires P;
  @ behavior <any name>:
  @   assumes \true;
  @   assigns L;
  @   ensures E;
  @*/
\end{listing}
Similarly, global \assigns\ and \ensures\ clauses are equivalent to a
single named behavior. More precisely, the following contract
\begin{listing}{1}
/*@ requires P;
  @ assigns L;
  @ ensures E;
  @ behavior b$_1$: `\dots`
  @ behavior b$_2$: `\dots`
  @ `\dots`
  @*/
\end{listing}
is equivalent to
\begin{listing}{1}
/*@ requires P;
  @ behavior <any name>:
  @   assumes \true;
  @   assigns L;
  @   ensures E;
  @ behavior b$_1$: `\dots`
  @ behavior b$_2$: `\dots`
  @ `\dots`
  @*/
\end{listing}

\begin{example}
\label{ex:bsearch}
In the following, \lstinline|bsearch(t,n,v)| searches for element \lstinline|v|
in array \lstinline|t| between indices \lstinline|0| and \lstinline|n-1|.

\listinginput{1}{bsearch.c}
%
The precondition requires array \lstinline|t| to be allocated at least
from indices \lstinline|0|
to \lstinline|n-1|. The two named behaviors
correspond respectively to the successful behavior and the failing
behavior.

Since the function is performing a binary search, it requires the
array \lstinline|t| to be sorted in increasing order: this is the purpose of
the predicate named \lstinline|t_is_sorted| in the \assumes\ clause of the
behavior named \lstinline|failure|.

See~\ref{sec:loop_annot} for a continuation of this example.
\end{example}

\begin{example}
  The following function illustrates the importance of different
  \assigns{} clauses for each behavior.

  \listinginput{1}{cond_assigns.c}
%
  Its contract means that it assigns values pointed to by
  \lstinline|p| or by \lstinline|q|,
  conditionally on the sign of~\lstinline|n|.
\end{example}


\subsubsection{Completeness of behaviors}
\label{sec:compl-behav}
In a contract with named behaviors, it is not required that
the disjunction of the \lstinline|A$_i$| is true, \emph{i.e.} it is not 
mandatory to provide a ``complete'' set of behaviors.
If such a condition is seeked, it is possible to add the following clause
to a contract:
\index{complete behaviors}
\begin{listing-nonumber}
/*@ `\dots`
  @ complete behaviors b$_1$,`\dots`,b$_n$;
  @*/
\end{listing-nonumber}
It specifies that the set of behaviors \lstinline|b$_1$,$\ldots$,b$_n$| 
is complete \emph{i.e.} that
\begin{listing-nonumber}
R ==> (A$_1$ || A$_2$ || ... || A$_n$)
\end{listing-nonumber}
holds, where \lstinline|R| is the precondition of the contract.
The simplified version of that clause
\begin{listing-nonumber}
/*@ `\dots`
  @ complete behaviors;
  @*/
\end{listing-nonumber}
means that all behaviors given in the contract should be taken into account.

Similarly, it is not required that two distinct behaviors are disjoint.
If desired, this can be specified with the following clause:
\index{disjoint behaviors}
\begin{listing-nonumber}
/*@ ...
  @ disjoint behaviors b$_1$,`\dots`,b$_n$;
  @*/
\end{listing-nonumber}
It means that the given behaviors are pairwise disjoint \emph{i.e.}
that, for all distinct $i$ and $j$,
\begin{listing-nonumber}
R ==> ! (A$_i$ && A$_j$)
\end{listing-nonumber}
holds.
The simplified version of that clause
\begin{listing-nonumber}
/*@ ...
  @ disjoint behaviors;
  @*/
\end{listing-nonumber}
means that all behaviors given in the contract should be taken into account.
Multiple \lstinline|complete| and \lstinline|disjoint| sets of behaviors can
be given for the same contract.

\oldremark{Abandonne}{

  2) on veut faire des behaviors differents suivant le post-state
     -> clause 'post assumes'

  exemples : (voir aussi annexe (Error Handling without exception))

  3) On pourrait avoir envie de specifier que les behavior sont
     mutuellement exclusifs et/ou couvrent tous les cas SELON
     LEURS CLAUSES "post assumes".

     ON NE TRAITE PAS CE CAS S'IL N'Y A PAS DE BESOIN, et il est
     alors interdit d'utiliser les clauses "complete behavior" et
     "disjoint behavior" sur des behavior contenant des "post assumes"

     la semantique serait:

    La condition de d'exclusion mutuelle des cas de behaviors
    est: pour toute paire de behavior $B_1$ et $B_2$ : postcondition
    supplementaire
\[
           not ((\old(assumes_1) et postassume_1)
                 et
                (\old(assumes_2) et postassume_2))
\]
    La condition de couverture des cas de behaviors est :

    la disjunction des $(\old(assumes_i) et postassume_i)$ est une postconditionn supplementaire

}


\oldremark{SD+Patrick}{
  Stephane confirme qu'AIRBUS utilise dans leur methodologie la notion
  d'assumes, mais pas celle de post-assumes.

  La possibilit� d'ecrire des assumes permettra de reduire la taille
  des postconditions actuelles et conduira a moins de repetitions des
  conditions de behaviors.
  Par contre, placer les actuelles preconditions dans des behaviors
  devrait conduire a plus de repetitions dans les clauses requires.
  Il est important pour AIRBUS de pouvoir eviter les repetitions de termes
  dans les proprietes.

  Ne serait-il pas possible de d�finir dans les contrats de fonctions
  des pr�dicats ou termes parlant des operandes de la fonction, sans
  passer par des fonctions logiques auxquelles on passerait en
  argument les operandes de la fonctions ?

  Ces predicats ou termes aurait une portee locale a la fonction.
}


% \subsubsection{Error cases}

% See Section~\ref{sec:errorhandling} for a discussion about
% behaviors related to error cases.

\subsection{Memory locations and sets of terms}
\label{sec:locations}\indextt{location}

% PB: Modified since \separated operates on pointers and not on locations
There are several places where one needs to describe a set of memory
locations: in \assigns clauses of function contracts, or
in \Loop \assigns clauses (see section~\ref{sec:loop_annot}).
A \emph{memory location} is then any set of terms denoting a set of
l-values\index{l-value}. Moreover, a location given as argument to
an \assigns{} clause must be a set of modifiable l-values, as described in
Section~\ref{sec:glossary}. 
More generally, we introduce syntactic constructs to denote \emph{sets of
  terms} which are also usefull for \separated predicate (see Section~\ref{sec:separated})

\oldremark{PB+VP}{
  pr�ciser le statut de \textbackslash{}result: est-ce une modifiable
  l-value ou pas? D'un c�t�, ce n'en est pas vraiment une en C, mais
  on aimerait peut-�tre pouvoir �crire des assigns un peu raffin�s
  comme \texttt{assigns \textbackslash{}result.a \textbackslash{}from i = i;}}

\begin{figure}
  \fbox{\begin{minipage}{0.97\textwidth}
      \input{loc_modern.bnf}
    \end{minipage}}
  \caption{Grammar for sets of terms}
\label{fig:gram:locations}
\end{figure}

The grammar for sets of terms is given in
Figure~\ref{fig:gram:locations}. The semantics is given below,
where $s$ denotes any \textsl{tset}.
\begin{itemize}
\item \lstinline|\empty|\indexttbs{empty} denotes the empty set.
\item a simple term denotes a singleton set.
\item \lstinline|s->id| denotes the set of \lstinline|x->id| for each
  \lstinline|x $\in$ s|.
\item \lstinline|s.id| denotes the set of \lstinline|x.id| for each
  \lstinline|x $\in$ s|.
\item \lstinline|*s| denotes the set of \lstinline|*x| for each
  \lstinline|x $\in$ s|.
\item \lstinline|&s| denotes the set of \lstinline|&x| for each
  \lstinline|x $\in$ s|.
\item \lstinline|s$_1$[s$_2$]| denotes the set of
\lstinline|x$_1$[x$_2$]| for each \lstinline|x$_1$ $\in$ s$_1$|
  and \lstinline|x$_2$ $\in$ s$_2$|.
\item \lstinline|t$_1$ .. t$_2$| denotes the set of integers between 
\lstinline|t$_1$| and \lstinline|t$_2$|, included. 
If \lstinline|t$_1$ $>$ t$_2$|, this the same as \lstinline+\empty+
\item \lstinline|\union(s$_1$,$\ldots$,s$_n$)|\indexttbs{union}
denotes the union 
of \lstinline|s$_1, $s$_2,\ldots $| and \lstinline|s$_n$|;
\item \lstinline|\inter(s$_1$,$\ldots$,s$_n$)|\indexttbs{inter}
denotes the intersection 
of \lstinline|s$_1, $s$_2,\ldots $| and \lstinline|s$_n$|;
\item \lstinline|s$_1$+s$_2$| denotes the set of
  \lstinline|x$_1$+x$_2$| for each \lstinline|x$_1$ $\in$ s$_1$|
  and \lstinline|x$_2$ $\in$ s$_2$|;
\item \lstinline|(s)| denotes the same set as \lstinline|s|;
\item \lstinline+{ s | b ; P }+ denotes set
  comprehension\index{comprehension},
  that is the union of the sets denoted by \lstinline|s| for each
  value \lstinline|b| of binders
  satisfying predicate \lstinline|P|
  (binders \lstinline|b| are bound in both \lstinline|s| and \lstinline|P|).
\end{itemize}

Note that \lstinline|assigns \nothing| is equivalent to
\lstinline|assigns \empty|; it is left for convenience.

\begin{example}
  The following function sets to 0 each cell of an array.

  \listinginput{1}{assigns_array.c}
%
  It is annotated with three equivalent \assigns{} clauses, each one
  specifying that only the set of cells
  \lstinline|{t[0],$\ldots$,t[n-1]}|
  is modified.
\end{example}

\begin{example}
  \label{ex:locations-list}
  The following function increments each value stored in a linked
  list.

  \listinginput{1}{assigns_list.c}
%
  The \assigns{} clause specifies that the set of modified memory
  locations is the set of fields \lstinline|q->hd| for each pointer
  \lstinline|q|
  reachable from \lstinline|p| following \lstinline|next| fields. See
  Section~\ref{sec:inductivepredicates} for details about the
  declaration of the predicate \texttt{reachable}.
\end{example}

\subsection{Default contracts, multiple contracts}
\label{sec:multiplecontracts}

A C function can be defined only once but declared several times.
It is allowed to annotate each of these declarations with contracts.
Those contracts are seen as a single contract with the union of the
\requires\ clauses and behaviors.

On the other hand, a function may have no contract at all, or a
contract with missing clauses. Missing \lstinline|requires| and
\lstinline|ensures| clauses default to \lstinline!\true!.
If no \lstinline|assigns| clause is given, it remains unspecified. If the
function under consideration has only a declaration but no body, then
it means that it potentially modifies ``everything'', hence in practice
it will be impossible to verify anything about programs calling that
function; in other words giving it a contract is in practice
mandatory. On the other hand, if that function has a body, giving no
\lstinline|assigns| clause means in practice that it is left to tools to
compute an over-approximation of the sets of assigned locations.


\oldremark{Claude}{

  Attention, cette semantique est dangereuse car si en pratique on
  inclut une partie seulement du contrat, on peut oublier de verifier
  un requires.

  Autres solutions:

  1) interdire les contrats multiples

  2) proposer une syntaxe pour distinguer les contrats initiaux, et
     les contrats complementaires (comme le also de JML)

}

\section{Statement annotations}
\index{annotation}
\oldremark{Claude}{
  Il faut absolument ajouter des regles de portee, a plusieurs
  endroits dans ce document

  Note c'est maintenant fait pour les \at
}

Annotations on C statements are of three kinds:
\begin{itemize}
\item Assertions: \index{assertion}allowed before any C statement or
  at end of blocks.
\item\index{annotation!loop}\index{loop!annotation}
  Loop annotations: \invariant, \assigns clause, \variant; allowed
  before any loop statement: \lstinline|while|, \lstinline|for|, and
  \lstinline|do ... while|.
\item\index{statement contract}\index{contract}
  Statement contracts: allowed before any C statement, specifying
  their behavior in a similar manner to C function contracts.
\end{itemize}

\subsection{Assertions}
\indextt{assert}
\label{sec:assertions}
\begin{figure}[t]
  \begin{cadre}
    \input{assertions_modern.bnf}
  \end{cadre}
  \caption{Grammar for assertions}
  \label{fig:gram:assertions}
\end{figure}

The syntax of assertions is given in Figure~\ref{fig:gram:assertions},
as an extension of the grammar of C statements.

\begin{itemize}
\item
  \lstinline|assert P|  means that \lstinline|P| must hold in the current state
  (the sequence point where the assertion occurs).

\item The variant \lstinline|for id$_1$,$\ldots$,id$_k$: assert P|
  associates the assertion to the named behaviors \lstinline|id$_i$|, each
  of them being a behavior identifier for the current function (or a
  behavior of an enclosing block as defined later in
  Section~\ref{sec:statement_contract}).  It means that this assertion
  must hold only for the considered behaviors.
\end{itemize}

\oldremark{Claude}{unifier les 2 derniers cas? car a-t-on besoin de
  nommer un assert en dehors d'un behavior de la fonction? Surtout que
  l'on peut deja ecrire assert (id : P)
}

\subsection{Loop annotations}
\label{sec:loop_annot}

\begin{figure}[t]
  \begin{cadre}
    \input{loops_modern.bnf}
  \end{cadre}
  \caption{Grammar for loop annotations}
  \label{fig:gram:loops}
\end{figure}

The syntax of loop annotations is given in Figure~\ref{fig:gram:loops},
as an extension of the grammar of C statements.
% [Patrick B.] START proposition:
\textsl{Loop-allocation} clauses allow to specify which memory location 
is dynamically allocated or deallocated by a loop from the \textsl{heap}; 
they are defined later in Section~\ref{sec:allocation-clause}.
% END proposition

\subsubsection{Loop invariants and loop assigns}
\index{invariant!loop}\index{loop!invariant}\index{loop!assigns}
The semantics of loop invariants and loop assigns is defined as follows: 
a simple loop annotation of the form
\begin{listing}{1}
/*@ loop invariant I;
  @ loop assigns L;
  @*/
...
\end{listing}
specifies that the following conditions hold.
\begin{itemize}
\item The predicate \lstinline|I| holds before entering the loop
  (in the case of a
  \lstinline|for| loop, this means right after the initialization expression).
\item The predicate \lstinline|I| is an inductive invariant, that is
  if \lstinline|I| is assumed true in some state where the condition
  \lstinline|c| is also true, and if execution of the loop body in
  that state ends normally at the end of the body or with a
  \Continue statement, \lstinline|I| is true in the
  resulting state. 
  If the loop condition has side effects, these are
  included in the loop body in a suitable way:
  \begin{itemize}
  \item for a \lstinline|while (c) s| loop, \lstinline|I| must be
    preserved by the side-effects of \lstinline|c| followed by
    \lstinline|s|;
  \item for a \lstinline|for(init;c;step) s| loop, \lstinline|I| must
    be preserved by the side-effects of \lstinline|c| followed by
    \lstinline|s| followed by \lstinline|step|;
  \item for a \lstinline|do s while (c);| loop, \lstinline|I| must be
    preserved by \lstinline|s| followed by the side-effects of
    \lstinline|c|.
  \end{itemize}
Note that if \lstinline|c| has side-effects, the invariant might not
be true at the exit of the loop: the last ``step'' starts from a state
where \lstinline|I| holds, performs the side-effects of \lstinline|c|,
which in the end evaluates to false and exits the loop. Likewise, if a
loop is exited through a \Break statement, \lstinline|I|
does not necessarily hold, as side effects may occur between
the last state in which \lstinline|I| was supposed to hold and 
the \Break statement.

\item At any loop iteration, any location that was allocated before
  entering the loop, and is not member of \lstinline|L| (interpreted in the
  current state) has the same value as
  before entering the loop.
  In fact, the \Loop \assigns clause specifies an inductive
  invariant for the locations that are not member of \lstinline|L|.
\end{itemize}

\subsubsection{Loop behaviors}
\index{loop!behavior}
A loop annotation preceded by
\lstinline|for id_1,$\ldots$,id_k:| is similar as above, but
applies only for behaviors \lstinline|id_1|,$\ldots$,\lstinline|id_k|
of the current function, 
hence in particular holds only under the assumption of their \assumes
clauses.


\paragraph{Remarks}
\begin{itemize}
\item The \old{} construct is not allowed
  in loop annotations. The \lstinline|\at| form should be used to
  refer to another state (see Section~\ref{sec:at}).
\item When a loop exits with \lstinline{break} or \lstinline{return} or
  \lstinline{goto}, it is not required that the loop invariant holds.
  In such cases, locations that are not member of $L$ can be
  assigned between the end of the previous iteration and the exit statement.
\end{itemize}

\begin{example}
\label{ex:bsearch2}
Here is a continuation of example~\ref{ex:bsearch}. Note the use of
a loop invariant associated to a function behavior.

\listinginput{1}{bsearch2.c}

\end{example}

\subsubsection{Loop variants}\label{sec:loop-variant}
\indextt{variant}\index{termination}\index{loop!variant}

Optionally, a loop annotation may include a loop variant of the form
\begin{listing-nonumber}
/*@ loop variant m; */
\end{listing-nonumber}
where \lstinline|m| is a term of type \lstinline|integer|.

The semantics is as follows: for each loop iteration that terminates
normally or with \lstinline|continue|, the value of \lstinline|m| at end of the
iteration must be smaller that its value at the beginning of the
iteration. Moreover, its value at the beginning of the iteration 
must be nonnegative. Note that the value of \lstinline|m| at loop exit might be
negative. It does not compromise termination of the
loop. Here is an example:

\begin{example}
  \hspace{0pt}
  \label{ex:loopvariant}
  \listinginput{1}{loopvariantnegative.c}
\end{example}

It is also possible to specify termination orderings other than the
usual order on integers, using the additional \lstinline|for|
modifier. This is explained in Section~\ref{sec:termination}.

\subsubsection{General inductive invariants}
\index{invariant}
It is actually allowed to pose an inductive invariant
anywhere inside a loop body.
For example, it makes sense for a
\lstinline|do s while (c);| loop to contain an invariant right after statement
\lstinline|s|. Such an invariant is a kind of assertion, as
shown in Figure~\ref{fig:advancedinvariants}.

\begin{figure}[t]
  \begin{cadre}
    \input{generalinvariants_modern.bnf}
  \end{cadre}
  \caption{Grammar for general inductive invariants}
  \label{fig:advancedinvariants}
\end{figure}

\begin{example}
  In the following example, the natural invariant holds at this point
  (\lstinline!\max! and \lstinline!\lambda! are introduced later in
  Section~\ref{sec:higherorder}).
  It would be less convenient to set an invariant at the beginning
  of the loop.
  \listinginput{1}{dowhile.c}
\end{example}

More generally, loops can be introduced by \lstinline|goto|s.
As a consequence, such invariants may occur anywhere
inside a function's body. The meaning is that the invariant
holds at that point, much like an \lstinline{assert}. Moreover, the
invariant must be inductive, \emph{i.e.} it must be preserved across a
loop iteration. Several invariants are allowed at different places in
a loop body. These extensions are useful when dealing with complex
control flows.

\oldremark{Claude}{TODO(citer une ref
a propos des invariants inductifs)
}

\begin{example}
Here is a program annotated with an invariant inside the loop body:
\listinginput{1}{advancedloopinvariants.c}
The control-flow graph of the code is as follows
\begin{center}
\includegraphics{cfg.mps}
%\begin{tikzpicture}
%\path (0,4) node[fill=yellow,draw](n1){1}
%      (0,2) node[fill=yellow,draw](n2){2};
%\draw[->,very thick] (n1.south) -- (n2.north);
% \filldraw[fill=green!20] ellipse (1cm and 4mm);
% \path (-2,2) node[fill=yellow,draw] (eb) {Expected behavior}
%       (-2,1) node[fill=yellow,draw] (fs) {Formal Specification}
%       (0,0) node (proof) {Proofs}
%       (2,1) node[fill=blue!30,draw] (pm) {Program Model}
%       (2,2) node[fill=blue!30,draw] (prog) {Program};
% \draw[->,very thick] (eb.south) -- (fs.north);
% \draw[->,very thick,shorten >=5mm] (fs.south) -- (proof.north);
% \draw[->,very thick] (prog.south) -- (pm.north);
% \draw[->,very thick,shorten >=5mm] (pm.south) -- (proof.north);
%\end{tikzpicture}
\end{center}
The invariant is inductively preserved by the two paths that go from
node ``inv'' to itself.
\end{example}

\begin{example}
The program
\begin{listing}{1}
  int x = 0;
  int y = 10;

  /*@ loop invariant 0 <= x < 11;
    @*/
  while (y > 0) {
    x++;
    y--;
  }
\end{listing}
is not correctly annotated, even if it is true that \lstinline|x|
remains smaller
than 11 during the execution. This is because it is not true that the
property \lstinline|x<11| is preserved by the execution of
\lstinline|x++ ; y--;|. A
correct loop invariant could be \lstinline|0 <= x < 11 && x+y == 10|. It holds
at loop entrance and is preserved (under the assumption of the loop
condition \lstinline|y>0|).

Similarly, the following general invariants are not inductive:
\begin{listing}{1}
  int x = 0;
  int y = 10;

  while (y > 0) {
    x++;
    //@ invariant 0 < x < 11;
    y--;
    //@ invariant 0 <= y < 10;
  }
\end{listing}
since \lstinline|0 <= y < 10| is not a consequence of
hypothesis \lstinline|0 < x < 11| after executing \verb|y--|; and
\lstinline|0 < x < 11| cannot be deduced from \lstinline|0 <= y < 10| after
looping back
through the condition \lstinline|y>0| and executing \lstinline|x++|. Correct
invariants could be:
\begin{listing}{1}
  while (y > 0) {
    x++;
    //@ invariant 0 < x < 11 && x+y == 11;
    y--;
    //@ invariant 0 <= y < 10 && x+y == 10;
  }
\end{listing}

\end{example}

\subsection{Built-in construct \texorpdfstring{\at}{\textbackslash{}at}}
\label{sec:at}
\indexttbs{at} Statement annotations usually need another additional
construct \lstinline|\at(e,id)| referring to the value of the
expression \lstinline|e| in the state at label \lstinline|id|. 
In particular, for a C array of \lstinline|int|, 
\lstinline|t|, \lstinline|\at(t,id)| is a logical
array whose content is the same as the one of \lstinline|t| in state at label
\lstinline|id|. It is thus very different from 
\lstinline|\at((int *)t,id)|, which
is a pointer to the first element of \lstinline|t| (and stays the same between 
the state at \lstinline|id| and the current state). Namely, if 
\lstinline|t[0]| has changed since \lstinline|id|, we have
\lstinline|\at(t,id)[0] != \at((int *)t,id)[0]|.

The label \lstinline|id| can be either a regular C label, or a label added
within a ghost statement as described in Section~\ref{sec:ghost}. This
label must be declared in the same function as the occurrence of
\lstinline|\at(e,id)|, but unlike \lstinline|goto|s, more restrictive
scoping rules must be respected:
\begin{itemize}
\item the label \lstinline|id| must occur before the occurrence of
\lstinline|\at(e,id)|
  in the source;
\item the label \lstinline|id| must not be inside an inner block.
\end{itemize}
These rules are exactly the same rules as for the visibility of local
variables within C statements (see~\cite{KR88}, Section A11.1).

\subsubsection*{Default logic labels}\label{sec:default-logic-labels}
\indextt{Pre}\indextt{Here}\indextt{Old}\indextt{Post}
There are four predefined logic labels: \lstinline|Pre|, \lstinline|Here|,
\lstinline|Old| and
\lstinline|Post|. \lstinline|\old(e)|\indexttbs{old} is
in fact syntactic sugar for \lstinline|\at(e,Old)|.

\begin{itemize}

\item The label \lstinline|Here| is visible in all statement annotations,
  where it refers to the state where the annotation appears; and in
  all contracts, where it refers to the pre-state for the
  \lstinline|requires|, \lstinline|assumes|, \lstinline|assigns|,
  \lstinline|variant|,
  \lstinline|terminates|,
  clauses and the post-state for other clauses.
It is also visible in data invariants, presented in Section~\ref{sec:invariants}.
\item The label \lstinline|Old| is visible in \lstinline|assigns| and
  \lstinline|ensures| clauses of all contracts (both for functions and for
  statement contracts described below in
  Section~\ref{sec:statement_contract}), and refers to the pre-state
  of this contract.
% [Patrick B.] START remarque:
% ``Old'' peut �galement s'utiliser dans les clauses relatives aux
% terminaisons abruptes des contrats d'instructions. Ce n'est pas dit ici.
% END remarque
\item The label \lstinline|Pre| is visible in all statement annotations,
  and refers to the pre-state of the function it occurs in.
\item The label \lstinline|Post| is visible in
\lstinline|assigns| and \lstinline|ensures| clauses of all contracts,
and it refers to the post-state.

\end{itemize}

There is one particular case for \lstinline|assigns| and \lstinline|ensures|
clauses of function contracts where formal parameters\index{formal parameter}
of functions cannot refer to the label \lstinline|Post|.
In such clauses formal parameters always refer implicitly to the label
\lstinline|Pre|, and any \lstinline|\at| construct can modify the
interpretation of formal parameters.

No logic label is visible in global logic declarations
such as lemmas, axioms, definition of predicate or logic
functions. When such an annotation needs to refer to a given memory
state, it has to be given a label binder: this is described in
Section~\ref{sec:logicalstates}.

\begin{example}
  The code below implements the famous extended Euclid's algorithm
  for computing the greatest common divisor of two integers $x$ and
  $y$, while computing at the same time two \emph{B\'ezout coefficients} $p$
  and $q$ such that $p\times x+q\times y = \gcd(x,y)$.
  The loop invariant for the B\'ezout property needs to refer to the
  value of $x$ and $y$ in the pre-state of the function.
  \listinginput{1}{euclide.c}
\end{example}

\begin{example}
  Here is a toy example illustrating tricky issues with \lstinline|\at| and
  labels:

  \listinginput{1}{oldat.c}

  The two \lstinline|ensures| clauses are equivalent. The simpler clause
  \lstinline|t[i] == \old(t[i]) + 1| would be wrong because in
  \lstinline|\old(t[i])|, \lstinline|i| denotes the value of $i$ n the
  pre-state.

  Also, the \lstinline|assigns| clause \lstinline|i,t[i]| would be
  wrong too because again in \lstinline|t[i]|, the value of $i$ in the
  pre-state is considered.

\end{example}

\subsection{Statement contracts}
\label{sec:statement_contract}
\index{statement contract}\index{contract}
\begin{figure}[t]
  \begin{cadre}
    \input{st_contracts_modern.bnf}
  \end{cadre}
  \caption{Grammar for statement contracts}
  \label{fig:gram:stcontracts}
\end{figure}

The grammar for statement contracts is given in
Figure~\ref{fig:gram:stcontracts}. It is similar to function
contracts, but without \decreases{} clause. Additionally, a statement contract
may refer to enclosing named behaviors, with the form
\lstinline|for id:...|.
Such contracts are only valid for the
corresponding behaviors, in particular only under the
corresponding \assumes{} clause.

\begin{todo}
  Add an example with a \For{}.
\end{todo}

% [Patrick B.] START proposition:
The \ensures\indextt{ensures} clause does not constrain the
post-state when the annotated statement is terminated 
by a \lstinline|goto| jumping out of it,
by a \Break{}, \Continue{} or \Return{} statement,
or by a call to the \exit{} function.
To specify such behaviors, \textsl{abrupt clauses} (described in
Section~\ref{sec:abrupt-clause}) need to be used.

On the other hand, it is different with \assigns\indextt{assigns} clauses.
The locations having their value modified during the path execution, starting
at the begining of the annoted statement and leading 
to a \lstinline|goto| jumping out of it, should be part of its 
\assigns clause.

\begin{example} The clause \lstinline|assigns \nothing;| does not hold for that statement,
even if the clause \lstinline|ensures x==\old(x);| holds: 
\begin{listing}{1}
/*@ assigns x;
  @ ensures x==\old(x);
  @*/
  if (c) {
    x++;
    goto L;
  }
L: ...
\end{listing}
\end{example}
% END proposition

% [Patrick B.] START proposition:
\textsl{Allocation-clauses} allow to specify which memory location 
is dynamically allocated or deallocated by the annotated statement from the \textsl{heap}; 
they are defined later in Section~\ref{sec:allocation-clause}.
% END proposition


\section{Termination}
\label{sec:termination}
\index{termination}
The property of termination concerns both loops and recursive function
calls.
% For that purpose, loops can be annotated with a \texttt{loop
%  variant}\index{loop variant} clause, and functions can be annotated
%with a \decreases{} clause.
Termination is guaranteed by attaching a measure function to each loop
(aspect already addressed in Section~\ref{sec:loop-variant}) 
and each recursive function.
By default, a measure is an
integer expression, and measures are compared using the usual ordering
over integers (Section~\ref{sec:integermeasures}). It is also possible
to define
measures into other domains and/or using a different ordering relation
(Section~\ref{sec:generalmeasures}).

\subsection{Integer measures}
\label{sec:integermeasures}
\indextt{decreases}\indextt{variant}
Functions are annotated with integer measures with the syntax
\begin{listing-nonumber}
//@ decreases e;
\end{listing-nonumber}
and loops are annotated similarly with the syntax
\begin{listing-nonumber}
//@ loop variant e;
\end{listing-nonumber}
where the logic expression \lstinline|e| has type
\lstinline|integer|.
For recursive calls, or for loops, this expression must decrease for
the
relation \lstinline|R| defined by
\begin{listing-nonumber}
R(x,y) <==> x > y && x >= 0.
\end{listing-nonumber}
In other words, the measure must be a decreasing sequence of integers
which remain nonnegative, except possibly for the last value of the
sequence (See example~\ref{ex:loopvariant}).

\begin{example}
  The clause \lstinline|loop variant u-l;| can be added to the loop 
  annotations of the example~\ref{ex:bsearch2}.
  The measure \lstinline|u-l| decreases
  at each iteration, and remains nonnegative, except at the last
  iteration where it may become negative.
\begin{listing}{16}
  @ ...
  @ loop variant u-l; */
while ...
\end{listing}
\end{example}

\subsection{General measures}
\label{sec:generalmeasures}

More general measures on other types can be provided, using the
keyword \lstinline|for|. For functions it becomes
\begin{listing-nonumber}
//@ decreases e for R;
\end{listing-nonumber}
and for loops
\begin{listing-nonumber}
//@ loop variant e for R;
\end{listing-nonumber}
In those cases, the logic expression \lstinline|e| has some type
$\tau$ and \lstinline|R|
must be relation on $\tau$, that is a binary predicate declared 
(see Section~\ref{sec:logicspec} for details) as
\begin{listing-nonumber}
//@ predicate R($\tau$ x, $\tau$ y) $\cdots$
\end{listing-nonumber}
Of course, to guarantee
termination, it must be proved that \lstinline|R| is a well-founded relation.

\begin{example}
  The following example illustrates a variant annotation
  using a pair of integers, ordered lexicographically.
  \listinginput{1}{lexico.c}
\end{example}

\subsection{Recursive function calls}

The precise semantics of measures on recursive calls, especially in
the general case of mutually recursive functions, is given as follows.
We call \emph{cluster} a set of mutually recursive functions which
is a strongly connected component of the call graph.
Within each cluster, each function must be annotated with a \decreases\
clause with the same relation \lstinline|R| (syntactically). Then, in the body
of any function \lstinline|f| of that cluster, any recursive call to a function
\lstinline|g| must occur in a state where the measure attached to
\lstinline|g| is smaller
(w.r.t \lstinline|R|) than the measure of \lstinline|f| in the pre-state of \lstinline|f|. This also
applies when \lstinline|g| is \lstinline|f| itself.

\begin{example}
  Here are the classical factorial and Fibonacci functions:
  \listinginput{1}{fact.c}
\end{example}

\begin{example}
  This example illustrates mutual recursion:
  \listinginput{1}{mutualrec.c}
\end{example}

% \subsection{Function termination}

% TODO: a mettre plutot dans les contrats

% clause 'terminates P' specifie que si P est vrai dans le pre-state,
% alors la fonction ou le statement objet du contrat se termine en temps
% fini.

% pas de valeur par defaut

% Patrick propose plutot la negation: 'diverges P' car ce que l'on veut
% specifier c'est que la fonction ne boucle indefiniment.

% en faveur de prendre 'diverges' : c'est comme ca en JML.


\subsection{Non-terminating functions}
\label{sec:non-term-funct}
\indextt{terminates}
\experimental

There are cases where a function is not supposed to terminate. For
instance, the \lstinline|main| function of a reactive program might be a
\lstinline|while(1)| which indefinitely waits for an event to process. More
generally, a function can be expected to terminate only if some
preconditions are met. In those cases, a \lstinline|terminates| clause can
be added to the contract of the function, under the following form:

\begin{listing-nonumber}
//@ terminates p;
\end{listing-nonumber}

The semantics of such a clause is as follows: if \lstinline|p| holds, then the
function is guaranteed to terminate (more precisely, its
termination must be proved). If such a clause is not present (and in
particular if there is no function contract at all), it
defaults to \lstinline|terminates \true;| that is the function is supposed
to always terminate, which is the expected behavior of most
functions.

Note that nothing is specified for the case where \lstinline|p| does not hold:
the function may terminate or not. In particular,
\lstinline|terminates \false;| does not imply that the function loops
forever. A possible specification for a function that never terminates
is the following:
\begin{listing}{1}
/*@ ensures \false;
    terminates \false;
*/
void f() { while(1); }
\end{listing}

\begin{example}
  A concrete example of a function that may not always terminate is
  the \lstinline|incr_list| function of example \ref{ex:locations-list}. In
  fact, The following contract is also acceptable for this function:
\listinginput{1}{terminates_list.c}
\end{example}

\section{Logic specifications}
\label{sec:logicspec}
\index{logic specification}\index{specification}
\begin{figure}[t]
  \fbox{\begin{minipage}{0.97\linewidth}\vfill \input{logic_modern.bnf}
    \vfill\end{minipage}}
  \caption{Grammar for global logic definitions}
\label{fig:gram:logic}
\end{figure}
The language of logic expressions used in annotations can be extended
by declarations of new logic types, and new constants, logic functions
and predicates. These declarations follows the classical setting of
\emph{algebraic specifications}.
The grammar for these declarations is given in Figure~\ref{fig:gram:logic}.

\subsection{Predicate and function definitions}
New functions and predicates can be \emph{defined} by explicit
expressions, given after an equal sign.
\begin{example}
  The following definitions \listinginput{1}{signdef.c} illustrates the
  definition of a new predicate \lstinline{is_positive} with an integer
  parameter, and a new logic function \lstinline{sign} with a real
  parameter returning an integer.
\end{example}

\subsection{Lemmas}
Lemmas are user-given propositions, a facility that might help theorem
provers to establish validity of ACSL specifications.

\begin{example}
  The following lemma
  \listinginput{1}{mean.c}
  is a useful hint for program like binary search.
\end{example}

Of course, a complete verification of an ACSL specification has to
provide a proof for each lemma.

\subsection{Inductive predicates}
\label{sec:inductivepredicates}
\index{inductive predicates}

A predicate may also be defined by an inductive definition. The
grammar for those style of definitions is given on
Figure~\ref{fig:gram:inductive}.
\index{inductive definitions}
\begin{figure}[t]
  \fbox{\begin{minipage}{0.97\linewidth} \input{inductive_modern.bnf}
    \end{minipage}}
  \caption{Grammar for inductive definitions}
\label{fig:gram:inductive}
\end{figure}

In general, an inductive definition of a predicate \lstinline|P| has the form
\begin{listing}{1}
/*@ inductive P(x$_1$,$\ldots$,x$_n$) {
  @   case c$_1$ : p$_1$;
...
  @   case c$_k$ : p$_k$;
  @ }
  @*/
\end{listing}
where each \lstinline|c$_i$| is an identifier and each \lstinline|p$_i$| 
is a proposition.

The semantics of such a definition is that \lstinline|P| is the least fixpoint
of the cases, i.e. is the smallest predicate (in the sense that it is
false the most often) satisfying the propositions \lstinline|p$_1,\ldots,$p$_k$|.
With this general form, the existence of a least fixpoint is not
guaranteed, so tools might enforce syntactic conditions on the form of
inductive definitions. A standard syntactic restriction could be to
allow only propositions \lstinline|p$_i$| of the form
\begin{listing-nonumber}
\forall y$_1$,`\dots`,y$_m$, h$_1$ ==> $\cdots$ ==> h$_l$ ==> P(t$_1$,`\dots`,t$_n$)
\end{listing-nonumber}
where \lstinline|P| occurs only positively in hypotheses \lstinline|h$_1,\ldots,$h$_l$|
(definite Horn clauses,
\url{http://en.wikipedia.org/wiki/Horn_clause}).

\begin{example}
  The following introduce a predicate \lstinline|isgcd(x,y,d)| meaning
  that \lstinline|d| is the greatest common divisor of \lstinline|x|
  and \lstinline|y|.
  \listinginput{1}{isgcd.c}
  This definition uses definite Horn clauses, hence is consistent.
\end{example}

Example~\ref{ex:locations-list} already introduced an inductive definition of
reachability in linked-lists, and was also bases on definite Horn
clauses thus consistent.

\subsection{Axiomatic definitions}

Instead of an explicit definition, one may introduce an
\emph{axiomatic} definitions for a set of types, predicates and logic
functions, which amounts to declare the expected profiles and a set of
axioms. The grammar for those constructions is given on
Figure~\ref{fig:gram:logicdecl}.

\index{axiomatic}
\begin{figure}[t]
  \fbox{\begin{minipage}{0.97\linewidth} \input{logicdecl_modern.bnf}
    \end{minipage}}
  \caption{Grammar for axiomatic declarations}
\label{fig:gram:logicdecl}
\end{figure}

\begin{example}
  The following axiomatization introduce a theory of finite lists of
  integers a la LISP.
  \listinginput{1}{intlists.c}
\end{example}

Like inductive definitions, there is no syntactic conditions which
would guarantee axiomatic definitions to be consistent. It is usually
up to the user to ensure that the introduction of axioms does not lead
to a logical inconsistency.

\begin{example}
  The following axiomatization
  \listinginput{1}{sign.c}
  is inconsistent since
  it implies \lstinline|sign(0.0) == 1| and \lstinline|sign(0.0) == -1|,
  hence \lstinline|-1 == 1|
\end{example}

\subsection{Polymorphic logic types}\label{sec:polym-logic-types}
\index{type!polymorphic}
\index{polymorphism}
\experimental

We consider here an algebraic specification setting based on
multi-sorted logic, where types can be \emph{polymorphic} that is
parametrized by other types. For example, one may declare the type of
polymorphic lists as
\listinginput{1}{listdecl.c}
One can then consider for
instance list of integers (\lstinline|list <integer>|), list of pointers
(e.g. \lstinline|list <char*>|), list of list of reals
(\lstinline|list<list <real> >|\footnote{In this latter case, note that the two
  '\texttt{>}' must be separated by a space, to avoid confusion with
  the shift operator.}), etc.

The grammar of Figure~\ref{fig:gram:logic} contains rules for
declaring polymorphic types and using polymorphic type
expressions.
%Notice that type variables are identifiers preceded by
%the quote character.


\subsection{Recursive logic definitions}
\index{recursion}
Explicit definitions of logic functions and predicates can be
recursive. Declarations in the same bunch of logic declarations are
implicitly mutually recursive, so that mutually recursive functions are
possible too.

\begin{example}
  The following logic declaration
  \listinginput{1}{max_index.c}
  defines a logic function which returns the maximal index \lstinline|i| between
  \lstinline|0| and \lstinline|n-1| such that \lstinline|t[i]=0|.
\end{example}

%\begin{example} The following introduce $n$-ary trees using list of children.
%  \input{mutualrec.pp}
%\end{example}

There is no syntactic condition on such recursive
definitions, such as limitation to primitive recursion. In essence, a
recursive definition of the form \lstinline+f(args) = e;+ where
\lstinline+f+ occurs in expression \lstinline+e+ is just a shortcut
for an axiomatic declaration of \lstinline+f+ with an axiom
\lstinline+\forall args; f(args) = e+.  In other words, recursive
definitions are not guaranteed to be consistent, in the same way that
axiomatics may introduce inconsistency. Of course, tools might provide
a way to check consistency.

\subsection{Higher-order logic constructions}
\label{sec:higherorder}

\experimental

\begin{figure}[t]
  \fbox{\begin{minipage}{0.98\linewidth}
      \input{higherorder_modern.bnf}
    \end{minipage}}
  \caption{Grammar for higher-order constructs}
\label{fig:gram:higherorder}
\end{figure}

Figure~\ref{fig:gram:higherorder} introduces new term constructs for
higher-order logic.
\begin{description}
\item[Abstraction]\indexttbs{lambda}
  The term \lstinline|\lambda $\tau_1$ x$_1$,$\ldots$,$\tau_n$ x$_n$; t|
  denotes the $n$-ary logic function which maps \lstinline|x$_1,\ldots,$x$_n$| to
  \lstinline|t|. It has the same precedence as \Forall and \Exists
\item[Extended quantifiers]
  Terms
 \lstinline|$\backslash quant$(t$_1$,t$_2$,t$_3$)|
  where $quant$ is \lstinline|max|\indexttbs{max},
  \lstinline|min|\indexttbs{min},
  \lstinline|sum|\indexttbs{sum},
  \lstinline|product|\indexttbs{product}
  or \lstinline|numof|\indexttbs{numof}
  are extended quantifications. \lstinline|t$_1$| and \lstinline|t$_2$| must have type
  \lstinline|integer|, and \lstinline|t$_3$| must be a unary function with an integer
  argument, and a numeric value (integer or real) except for
  \lstinline|\numof| for which it should have a boolean value. Their
  meanings are given as follows:
  \begin{tabular}{rcl}
    \lstinline|\max(i,j,f)| &=&
    $\max \{ \lstinline|f(i)|,\lstinline|f(i+1)|, \ldots, \lstinline|f(j)| \}$\\
    \lstinline|\min(i,j,f)| &=&
    $\min \{ \lstinline|f(i)|,\lstinline|f(i+1)|, \ldots, \lstinline|f(j)| \}$\\
    \lstinline|\sum(i,j,f)| &=&
    $\lstinline|f(i)|+\lstinline|f(i+1)|+\cdots+\lstinline|f(j)|$ \\
    \lstinline|\product(i,j,f)| &=&
    $\lstinline|f(i)|\times\lstinline|f(i+1)|\times\cdots\times\lstinline|f(j)|$
    \\
    \lstinline|\numof(i,j,f)| &=&
    $\#\{ \lstinline|k| \mid \lstinline|i| \leq \lstinline|k| \leq
          \lstinline|j| \land \lstinline|f(k)| \}$ \\
&=& \lstinline|\sum(i,j,\lambda integer k ; f(k) ? 1 : 0)|
  \end{tabular}

  If \lstinline|i>j| then \lstinline|\sum| and \lstinline|\numof| above are 0,
  \lstinline|\product| is 1, and \lstinline|\max| and \lstinline|\min| are
  unspecified (see Section~\ref{sec:twovaluedlogic}).
\end{description}


\begin{example}
  \label{ex:higherorder}
  Function that sums the element of an array of doubles.
  \listinginput{1}{sum.c}
\end{example}

\subsection{Concrete logic types}\label{sec:concrete-logic-types}
\index{type!concrete}
\experimental

\begin{figure}[t]
  \fbox{\begin{minipage}{0.98\linewidth}
      \index{record}\index{sum}
      \input{logictypedecl_modern.bnf}
    \end{minipage}}
  \caption{Grammar for concrete logic types and pattern-matching}
\label{fig:gram:logictype}
\end{figure}

\oldremark{SD6+All}{

  faire une analogie avec les types sommes (de Caml, des spec algebriques)

  c'est plutot des types avec constructeurs dans les spec algebriques


}


Logic types may not only be declared but also be given a
definition. Defined logic types can be either record
types\index{record}\index{type!record}, or sum
types\index{sum}\index{type!sum}. These definitions may be recursive.
For record types, the field access notation $t.id$ can be used, and
for sum types, a pattern-matching construction is available.
Grammar rules for these additional constructions are given in
Figure~\ref{fig:gram:logictype}

\oldremark{All}{

TODO: ajouter les types tableaux logiques: tableaux fonctionnels.

sucre syntaxique: t[i] pour select(t,i),  et dans le cas de code ghost, t[i] = e pour t = store(t,i,e)

Example: n-reines.


}

\begin{example}
  The declaration
  \listinginput{1}{listdef.c}
  introduces a concrete definition of finite lists. The logic definition
  \listinginput{1}{listlengthdef.c}
  defines the length of a list by recursion and pattern-matching.
\end{example}


\subsection{Hybrid functions and predicates}
\label{sec:logicalstates}
\index{hybrid!function}
\index{hybrid!predicate}

Logic functions and predicates may take both (pure) C types and logic
types arguments. Such an hybrid predicate (or function)
can either be defined with the same syntax as before (or axiomatized).

Be it defined either directly by an expression or through a set of axioms,
an hybrid function (or predicate) usually depends on one or
more program points, because it depends upon memory states, \emph{via}
expressions such as:
\begin{itemize}
\item pointer dereferencing: \lstinline|*p|, \lstinline|p->f|;
\item array access: \lstinline|t[i]|;
\item address-of operator: \lstinline|&x|;
\item built-in predicate depending on memory: \valid
% \item others ?
\end{itemize}
To make such a definition safe, it is mandatory to add after the
declared identifier a set of labels, between curly braces, as shown on
Figure~\ref{fig:gram:logiclabels}. Expressions
as above must then be enclosed into the \at{} construct to refer to a
given label. However, to ease reading of such logic expressions, it
is allowed to omit a label whenever there is only one label in the
context.

\begin{figure}[t]
  \fbox{\begin{minipage}{0.97\linewidth}
      \input{logiclabels_modern.bnf}
    \end{minipage}}
  \caption{Grammar for logic declarations with labels}
\label{fig:gram:logiclabels}
\end{figure}

\begin{example}\label{ex:nb_occ}
  The following annotations declare a function which returns the
  number of occurrences of a given double in a memory block storing doubles
  between the given indexes, together with the related axioms.  It
  should be noted that without labels, this axiomatization would be
  inconsistent, since the function would not depend on the values
  stored in t, hence the two last axioms would say both that \lstinline|a==b+1|
  and \lstinline|a==b| for some \lstinline|a| and \lstinline|b|.
  \listinginput{1}{nb_occ.c}
\end{example}

\begin{example}
  This second example defines a predicate which indicates whether two
  memory blocks of the same size are a permutation of each other. It
  illustrates the use of more than a single label. Thus, the \lstinline|\at|
  operator is mandatory here. Indeed the two blocks
  may come from two distinct memory states. Typically, one of the post
  condition of a sorting function would be \lstinline|permut{Pre,Post}(t,t)|.
  \listinginput{1}{permut.c}
\end{example}

% resurected reads clause

\subsection{Memory footprint specification: \texorpdfstring{\lstinline|reads|}{reads} clause}

\begin{figure}[t]
  \fbox{\begin{minipage}{0.97\linewidth}
      \input{logicreads_modern.bnf}
    \end{minipage}}
  \caption{Grammar for logic declarations with \lstinline|reads| clauses}
\label{fig:gram:logicreads}
\end{figure}


\experimental

Logic declaration can be augmented with a \reads{} clause, with the
syntax given in Figure~\ref{fig:gram:logicreads}, which extends the
one of Figure~\ref{fig:gram:logic}. This feature allows to specify the
\emph{footprint} of an hybrid predicate or function, that is the set of
memory locations which it depends on. From such an information, one
might deduce properties of the form $f\{L_1\}(args) = f\{L_2\}(args)$
if it is known that between states $L_1$ and $L_2$, the memory changes are
disjoint from the declared footprint.
%is useful when a
%function or a predicate is not easily definable, but can be more
%easily axiomatized.

% VP: exemple redondant avec le suivant.
% \begin{example}
%   The following annotations declare a function which returns the number of
%   positive elements in an array of doubles between the given indexes
%   and the related axioms.
%   \input{num_of_pos.pp}
%   It should be noted that without the reads
%   clauses, this axiomatization would be inconsistent, since the
%   function would not depend on the values stored in t, hence the two
%   last axioms would say both that $a=b+1$ and $a=b$ for some $a$ and
%   $b$.

%   \oldremark{Claude}{
%     expliquer mieux

%     De plus, si l'on oublie la clause reads �a reste
%     correct, donc il faudrait prevoir un warning pour un predicat avec un
%     type pointeur en argument sans reads...
%   }
% \end{example}

\begin{remark}{Claude}
on peut indiquer des conditions semantiques de validite des axiomes:

pour tout etat L1, L2 qui sont identiques sur le footprint, P{L1} =>
P{L2} ou P est la formule de l'axiome

\end{remark}

\begin{example}
  The following is the same as example~\ref{ex:nb_occ} augmented with
  a \lstinline|reads| clause.
  \listinginput{1}{nb_occ_reads.c}
  If for example a
  piece of code between labels \lstinline|L_1| and \lstinline|L_2|
  modifies \lstinline|t[k]| for
  some index \lstinline|k| outside \lstinline|i..j|, then one can deduce that
  \lstinline|nb_occ{L_1}(t,i,j,e)==nb_occ{L_2}(t,i,j,e)|.
\end{example}

% \begin{example}
%   This second example defines a predicate which indicates whether two
%   arrays of the same size are a permutation of each other. It
%   illustrates the use of more than a single label. Thus, the \lstinline|\at|
%   operator is mandatory here. Indeed the two arrays
%   may come from two distinct memory states. Typically, one of the post
%   condition of a sorting function would be \lstinline+permut{Pre,Here}(t,t)+.
%   \listinginput{1}{permut_reads.c}
% \end{example}

% \oldremark{SD7+All}{

%   As shown above, hybrid predicates and functions declared with reads
%   clauses should be used with caution, since they do not depend only
%   on the values of their parameters but also in an implicit way on the
%   values stored in memory.

%   reutiliser l'exemple en question ? (en le faisant mieux)

% }

\subsection{Specification Modules}
\label{sec:specmodules}
\index{module}

\notimplemented{Specification modules can be provided} to encapsulate
several logic definitions, for example
\listinginput{1}{listmodule.c}
Module components are then accessible using a qualified notation like
\begin{notimplementedenv}\lstinline|List::length|\end{notimplementedenv}.

Predefined algebraic specifications can be provided as
libraries\index{library} (see
section~\ref{chap:lib}), and imported using a construct like
\listinginput{1}{import.c} where the file \lstinline|list.acsl| contains logic
definitions, like the \lstinline|List| module above.



\section{Pointers and physical adressing}
\label{sec:pointers} 
% [Patrick B.] START proposition:
Grammar for terms and predicates is
extended with new constructs given in Figure~\ref{fig:gram:memory}. 
\textsl{Location-address} denotes the address of a memory location. 
It is a set of terms of some pointer type as defined in Section~\ref{sec:locations}.

\begin{figure}[t]
  \fbox{\begin{minipage}{0.97\linewidth}
      \input{memory_modern.bnf}
    \end{minipage}}
  \caption{Grammar extension of terms and predicates about memory}
\label{fig:gram:memory}
\end{figure}
% END proposition:



\subsection{Memory blocks and pointer dereferencing}
\label{subsec:memory}
C memory is structured into allocated blocks that can come either from a 
declarator or a call to one of the \lstinline|calloc|, \lstinline|malloc| or
\lstinline|realloc| functions. Blocks are characterized by their base address, 
{\it i.e.} the address of the declared object (the first declared object 
in case of an array declarator), or the pointer returned by the allocating 
function (when the allocation succeed) and their length.

ACSL provides the following built-in functions to deal with allocated blocks.
All of them takes an optional label identifier as argument. 
The default value of that label is defined in Section~\ref{sec:default-logic-labels}.
\begin{itemize}
\item \lstinline|\base_addr{L}(p)|\indexttbs{base\_addr}
returns the base address of the allocated block
containing, at the label \lstinline|L|, the pointer \lstinline|p|
\\ \makebox[5mm]{} \lstinline|\base_addr{id} : void* $\ra$ char*|

\item \lstinline|\block_length{L}(p)|\indexttbs{block\_length}
  returns the length (in bytes) of the allocated block containing, 
  at the label \lstinline|L|, its argument pointer.
\\ \makebox[5mm]{} \lstinline|\block_length{id} : void* $\ra$ size_t|
\end{itemize}

In addition, dereferencing a pointer may lead to run-time errors. A pointer
\lstinline|p| is said to be \emph{valid} if \lstinline|*p| is guaranteed to
produce a definite value according to the C standard~\cite{standardc99}. The
following built-in predicates deal with this notion:
\begin{itemize}
\item \valid\indexttbs{valid} applies to a set of terms 
(see Section~\ref{sec:locations}) of some pointer type. 
\lstinline|\valid{L}(s)| holds if and only if
 dereferencing any $p\in \lstinline|s|$ is safe  at label \lstinline|L|, both for reading from 
\lstinline|*p| and writing to it. In particular,
\lstinline|\valid{L}(\empty)| holds for any label \lstinline|L|. 
\\ \makebox[5mm]{} \lstinline|\valid{id} : set<$\alpha$ *> $\ra$ boolean|
\item 
\lstinline|\valid_read|\indexttbs{valid_read}
applies to a set of terms of some pointer type and holds if and only
if it is safe to read from all the pointers in the set
\\ \makebox[5mm]{} \lstinline|\valid_read{id} : set<$\alpha$ *> $\ra$ boolean|
\end{itemize}

\lstinline|\valid{L}(s)| implies \lstinline|\valid_read{L}(s)| but the reverse is 
not true. In particular, it is allowed to read from a string literal, but not 
to write in it (see \cite{standardc99}, 6.4.5\S6).

The status of \valid and \lstinline|\valid_read| constructs depends on the
type of their argument. Namely, \lstinline|\valid{L}((int *) p)| and 
\lstinline|\valid{L}((char *)p)| are not equivalent. On the other hand, if we 
ignore potential alignment constraints, the following equivalence is true for any pointer \lstinline|p|:
\begin{listing-nonumber}
\valid{L}(p) <==> \valid{L}(((char *)p)+(0..sizeof(*p)))
\end{listing-nonumber}
and
\begin{listing-nonumber}
\valid_read{L}(p) <==> \valid_read{L}(((char *)p)+(0..sizeof(*p)))
\end{listing-nonumber}

Some shortcuts are provided:
\begin{itemize}
\item \lstinline|\null|\indexttbs{null} is an
  extra notation for the
  null pointer (\emph{i.e.} a shortcut for \lstinline|(void*)0|). 
  As in C itself (see \cite{standardc99}, 6.3.2.3\S3),
  the constant \lstinline|0| can have any pointer type. In addition, 
  \lstinline|\valid{L}(\null)| is always false, for any
  logic label \lstinline|L|.
  Of course, \lstinline|\valid_read{L}(\null)| is always false too.
% \footnote{or to \texttt{NULL}, if
%    annotations are pre-processed are the appropriate inclusion of
%    \texttt{stddef.h} is added}
\item \lstinline|\offset{L}(p)|%
  \indexttbs{offset}
  returns the offset between \lstinline|p| and its base address
\\ \makebox[5mm]{} \begin{tabular}{lll}
    \lstinline|\offset{id}| &:& \lstinline|void* $\ra$ size_t| \\
    \lstinline|\offset{L}(p)| &=& \lstinline|(char*)p - \base_addr{L}(p)| \\
  \end{tabular} \\
Again, if there is no alignment constraints,
  the following property holds: for any set of pointers
\lstinline|s| and label \lstinline|L|, \lstinline|\valid_read{L}(s)| if and only if for all \lstinline|p$\in$s|:
  \begin{listing-nonumber}
    \offset{L}(p) >= 0 && \offset{L}(p) + sizeof(*p) <= \block_length{L}(p)
  \end{listing-nonumber}

\end{itemize}

\subsection{Separation}\label{sec:separated} 

ACSL provides a built-in function to deal with separation of locations:

\begin{itemize}
% [Patrick B.] START proposition:
\item \lstinline|\separated|\indexttbs{separated}
applies to sets of terms (see Section~\ref{sec:locations}) of some pointer type. 
\lstinline|\separated(s$_1$,s$_2$)|
holds for any set of pointers \lstinline|s$_1$| and \lstinline|s$_2$| 
if and only if for all \lstinline|p$\in$s$_1$| and \lstinline|q$\in$s$_2$|:
\begin{listing-nonumber}
  forall integer i,j; 0 <= i < sizeof(*p), 0 <= j < sizeof(*q)
     ==>  (char*)p + i != (char*)q + j 
\end{listing-nonumber}

In fact, \lstinline|\separated| is an $n$-ary predicate.\\
\lstinline|\separated(s$_1$,..,s$_n$)|\indexttbs{separated} means
that for each
$i\neq j$, \lstinline|\separated(s$_i$,s$_j$)|.
% END proposition:
\end{itemize}


\subsection{Dynamic allocation and deallocation}\label{sec:allocation-clause}
\indextt{allocates}\indextt{frees}
\index{allocation}\index{deallocation}\index{dynamic allocation}

\experimental

\textsl{Allocation-clauses} allow to specify which memory location 
is dynamically allocated or deallocated.
The grammar for those constructions is given on Figure~\ref{fig:gram:allocation}.

\lstinline|allocates \nothing| and \lstinline|frees \nothing| are respectively 
equivalent to \lstinline|allocates \empty| and \lstinline|frees \empty|; 
it is left for convenience like for \assigns clauses.

\begin{figure}[t]
  \fbox{\begin{minipage}{0.97\linewidth}
      \input{allocation_modern.bnf}
    \end{minipage}}
  \caption{Grammar for dynamic allocations and deallocations}
\label{fig:gram:allocation}
\end{figure}

\subsubsection{Allocation clauses for function and statement contracts}

Clauses \allocates and \frees are tied together. The simple contract  
\begin{listing-nonumber}
/*@ frees P$_1$,P$_2$,`\dots`;
  @ allocates Q$_1$,Q$_2$,`\dots`;
  @*/
\end{listing-nonumber}
means that any memory address that do not belong to the union of sets of terms
\lstinline|P$_i$| and \lstinline|Q$_j$| has the same {\sl allocation status} 
(see below) in the post-state than in the pre-state. The only one difference 
between \lstinline|allocates| and \lstinline|frees| is that sets 
\lstinline|P$_i$| are evaluated 
in the pre-state, and sets \lstinline|Q$_i$| are evaluated in the post-state.

\allocates \lstinline|Q$_1$,$\ldots$,Q$_n$| is equivalent to the postcondition
\begin{listing-nonumber}
\forall char* p; 
\separated(\union(Q$_1$,$\ldots$,Q$_n$),p)==>(\allocation{Here}(p)==\allocation{Pre}(p)
                               && \base_addr{Here}(p)==\base_addr{Pre}(p)
                               && \block_length{Here}(p)==\block_length{Pre}(p)
                               && \valid{Here}(p)<==>\valid{Pre}(p)
                               && \valid_read{Here}(p)<==>\valid_read{Pre}(p))
\end{listing-nonumber}

In fact, like the \assigns clause does not specify which memory location is assigned,
the \textsl{allocation-clauses} do not specify which memory location 
is dynamically allocated or deallocated.
Pre-conditions and post-conditions should be added to complete specifications
about allocations and deallocations.
The following built-in constructs can be used for that:

\begin{itemize}
\item  \begin{notimplementedenv}\lstinline|\allocation{L}(p)|\end{notimplementedenv}%
  \indexttbs{allocation}
returns the allocation status of the block containing, at the label \lstinline|L|, the pointer \lstinline|p| 
\\ \makebox[5mm]{} \lstinline|\allocation{id} : void* $\ra$| 
\begin{notimplementedenv}\lstinline|{\static, \register, \automatic, \dynamic, \unallocated}|
\end{notimplementedenv}%
This function is such that for any pointer \lstinline|p| and label \lstinline|L|
\begin{listing-nonumber}
\allocation{L}(p) == \allocation{L}(\base_addr(p))
\end{listing-nonumber}
and
\begin{listing-nonumber}
\allocation{L}(p)==\unallocated ==> !\valid_read{L}(p)
\end{listing-nonumber}
\end{itemize}

Some usefull shortcuts are provided:

\begin{itemize}
\item \lstinline|\allocable{L}(p)|\indexttbs{allocable}
holds if and only if
the pointer \lstinline|p| refers, at the label \lstinline|L|, 
to the base address of 
an unallocated memory block.
\\ \makebox[5mm]{} \lstinline|\allocable{id} : void* $\ra$ boolean|
\\
For any pointer \lstinline|p| and label \lstinline|L|
\begin{listing-nonumber}
\allocable{L}(p) <==> (\allocation{L}(p)==\unallocated && p==base_addr{L}(p)).
\end{listing-nonumber}
  
\item \lstinline|\freeable{L}(p)|\indexttbs{freeable}
holds if and only if
the pointer \lstinline|p| refers, at the label \lstinline|L|, 
to an allocated memory block 
that can be safely released using the C function \lstinline|free|.
\\ \makebox[5mm]{} \lstinline|\freeable{id} : void* $\ra$ boolean|
\\
For any pointer \lstinline|p| and label \lstinline|L|
\begin{listing-nonumber}
\freeable{L}(p) <==> (\allocation{L}(p)==\dynamic && p==base_addr{L}(p)).
\end{listing-nonumber}
  
\item \lstinline|\fresh{L$_0$,L$_1$}(p,n)|
\indexttbs{fresh},
  indicates that \lstinline|p| refers to an allocated memory block at label 
  \lstinline|L$_1$|, 
  but that it is not the case at label \lstinline|L$_0$|. 
  The predicate ensures also that, at label \lstinline|L$_1$|, the length
  (in bytes) of the block allocated dynamically equals to \lstinline|n|.
\\ \makebox[5mm]{} \lstinline|\fresh{id,id} : void*, \integer $\ra$ boolean|
\\
For any pointer \lstinline|p| and labels \lstinline|L$_0$| and \lstinline|L$_1$|
\begin{listing-nonumber}
\fresh{L$_0$,L$_1$}(p,n) <==> (\allocable{L$_0$}(p) && 
                         \freeable{L$_1$}(p) && 
                         \block_length{L$_1$}(p)==n &&
                         \valid{L$_1$}((char*)p+(0..n))
\end{listing-nonumber}
\end{itemize}

\begin{example}
  \lstinline|malloc| and \lstinline|free| functions can be specified as follows.
  \listinginput{1}{malloc_free_fn.c}
  Default labels for constructs dedicated to memory are 
  such that logic label \lstinline|Here| can be omitted.
\end{example}

When a behavior has only one of the two allocation clauses,
the given clause specifies the whole set of memory addresses to consider.
This means that the set value for the other clause of that behavior defaults to \nothing.
Now, when none of the two allocation clauses is given, the meaning is different
for anonymous behaviors and named behaviors:
\begin{itemize}
\item a named behavior without allocation clause does not specify anything about
      allocations and deallocations. 
      The allocated and deallocated memory blocks are in fact specified by the 
      anonymous behavior of the contract.
      There is no condition to verify for these named behaviors about allocations 
      and deallocations;
\item for an anonymous behavior, no allocation clause means that there is no newly 
      allocated nor deallocated memory block. 
      That is equivalent to give \allocates \nothing.
\end{itemize}
These rules are such that contracts without any allocation clause should be considered 
as having only one \allocates \nothing leading to a condition to verify for 
each anonymous behavior.

\begin{example}
  More precise specifications can be given using named behaviors under the 
  assumption of \assumes clauses.
  \listinginput{1}{malloc-free2-fn.c}
  The behaviors named \lstinline|allocation| and \lstinline|deallocation| do not 
  need an allocation clause.
  For example, the allocation constraint of the \lstinline|allocation| behavior 
  is given by the clause \allocates \result of the anonymous behavior of the  
  \lstinline|malloc| function contract. 
  To set a stronger contraint into the behavior named \lstinline|no_allocation|,
  the clause \allocates \nothing should be given.
\end{example}

\subsubsection{Allocation clauses for loop annotations}
\index{loop!allocation}\index{loop!deallocation}

Loop annotations are complemented by similar clauses allowing to specify which 
memory location 
is dynamically allocated or deallocated by a loop.
The grammar for those constructions is given on Figure~\ref{fig:gram:allocation}.

\Loop \allocates and \Loop \frees have the same semantics.
\begin{listing-nonumber}
//@ loop allocates Q$_1$,Q$_2$,`\dots`;
\end{listing-nonumber}
means that any memory address that do not belong to the union of sets of 
terms
\lstinline|Q$_1$ $\cup$ Q$_2$ $\cup \dots$|, interpreted in the current state, 
has the same allocation status in the current state than 
before entering the loop. Namely, as for \Loop \assigns, 
\Loop \allocates defines a loop invariant.

More precisely, with \lstinline|L$_0$| labelling the state before entering
the loop, the clause \Loop \allocates \lstinline|Q$_1$,$\ldots$,Q$_n$|
is equivalent to the postcondition
\begin{listing-nonumber}
\forall char* p; 
\separated(\union(Q$_1$,$\ldots$,Q$_n$),p)==>(\allocation{Here}(p)==\allocation{L$_0$}(p)
                               && \base_addr{Here}(p)==\base_addr{L$_0$}(p)
                               && \block_length{Here}(p)==\block_length{L$_0$}(p)
                               && \valid{Here}(p)<==>\valid{L$_0$}(p)
                               && \valid_read{Here}(p)<==>\valid_read{L$_0$}(p))
\end{listing-nonumber}


\begin{example}
 \listinginput{1}{loop-frees.c}
The addresses of locations \lstinline|q[0..n]| are not modified by the loop, 
but their values are.
It is the reason why the \lstinline|loop frees| clause requires 
the use of the ghost label \lstinline|L0|,
in order to catch the set of the released memory blocks.
\end{example}

A {\sl loop-clause} without allocation clause implicitly contents \Loop \allocates \nothing.
That means the {\sl allocation status} is not modified by the loop body.
A {\sl loop-behavior} without allocation clause means that  
the allocated and deallocated memory blocks are in fact specified by the allocation clauses
of the {\sl loop-clauses}
(Grammar of {\sl loop-clauses} and {\sl loop-behaviors} is given in Figure~\ref{fig:gram:loops}).

\section{Sets as first-class values}\label{sec:sets}
\index{location}
Sets of terms, as defined in Section~\ref{sec:locations},
can be used as first-class values in
annotations. All the elements of such a set must share the same type
(modulo the usual implicit conversions).
Sets have the built-in type \lstinline|set<A>|\index{set@\texttt{set} type}
where \lstinline|A| is the type of terms contained in the set.

In addition, it is possible to consider sets of pointers to values of
different types. In this case, the set is of type
\lstinline|set<char*>| and
each of its elements \lstinline|e| is converted to
\lstinline|(char*)e + (0..sizeof(*e)-1)|.

\begin{example}
  The following example defines the \emph{footprint} of a
  structure, that is the set of locations that can be accessed from an
  object of this type.

  \listinginput{1}{footprint.c}

  In the first definition, since union is made with a
  \lstinline|set<char*>| and a \lstinline|set<int*>|,
  the result is a \lstinline|set<char*>| (accordingly to typing of union). In
  other words, the two definitions above are equivalent.

  This logic function can be used as argument of \separated or
  \assigns clause.
\end{example}

Thus, the \separated predicate satisfies the following property (with
\lstinline|s$_1$| of type \lstinline|set<$\tau_1$*>| and \lstinline|s$_2$| of type
\lstinline|set<$\tau_2$*>|)
\begin{listing}{1}
\separated(s$_1$,s$_2$) <==>
  (\forall $\tau_1$* p; \forall $\tau_2$* q;
     \subset(p,s$_1$) && \subset(q,s$_2$) ==>
       (\forall \integer i,j;
         0 <= i < \sizeof($\tau_1$) && 0 <= j < \sizeof($\tau_2$) ==>
           (char*)p + i != (char*)q + j))
\end{listing}

and a clause \assigns \lstinline|L$_1$,$\ldots$,L$_n$| is equivalent
to the postcondition
\begin{listing-nonumber}
\forall char* p; \separated(\union(&L$_1$,$\ldots$,&L$_n$),p) ==> *p == \old(*p)
\end{listing-nonumber}

\section{Abrupt termination}
\label{sec:abrupt-clause}
\index{abrupt clause}

\experimental

\begin{figure}[t]
  \fbox{\begin{minipage}{0.98\textwidth}
     \input{exitbehavior_modern.bnf}
    \end{minipage}}
    \caption{Grammar of contracts about abrupt terminations}
  \label{fig:gram:abrupt-clauses}
\end{figure}

The \ensures\indextt{ensures} clause of function and statement contracts
does not constrain the post-state when the annotated function and
statement terminates respectively abruptly. In such cases,
\textsl{abrupt clauses} can be used inside \textsl{simple clause} or
\textsl{behavior body}. The allowed constructs are shown in
Figure~\ref{fig:gram:abrupt-clauses}.

The clauses \lstinline|breaks|\indextt{breaks}, \continues{}\indextt{continues} 
and \returns{}\indextt{returns}
can only be found in a statement contract and
state properties on the program state which hold when the
annotated statement terminates abruptly with the corresponding
statement (\lstinline|break|, \Continue{} or \lstinline|return|).

Inside these clauses, the construct \lstinline|\old(e)|\indexttbs{old}
is allowed and denotes, like for statement contracts
\lstinline|ensures|, \assigns and \lstinline|allocates|, the value of
\lstinline|e| in the pre-state of the statement.
 More generally, the visibility in \textsl{abrupt clauses} of predefined 
logics labels\indextt{Pre}\indextt{Here}\indextt{Old}\indextt{Post}
 (presented in Section~\ref{sec:default-logic-labels}) is the 
same as in \ensures{} clauses.

For the \lstinline|returns| case, the \result\indexttbs{result}
construct is allowed (if the function is not returning \void) and is bound
to the returned value.

\begin{example}
  The following example illustrates each abrupt
  clause of a statement contracts.
  \listinginput{1}{abrupt_termination.c}
\end{example}

The \exits{} clause can be used in both function and statement
contracts to give behavioral properties to the \main{} function
or to any function that may exit the program,
\emph{e.g.} by calling the \exit{} function.

In such clauses, \lstinline|\old(e)|\indexttbs{old}
is allowed and denotes the value of \lstinline|e| in the pre-state of
the function or statement,
and \lstinline|\exit_status| is bound to the return code,
\emph{e.g} the value returned by \main{} or the argument passed to \exit{}.
The  construct \lstinline|\exit_status|\indexttbs{exit_status} can be used only in \exits{},
 \assigns{} and \allocates{} clauses. On the contrary,
\result cannot be used in \exits{} clauses.

\begin{example}
\label{ex:assigns-and-abrupt-termination}
  Here is a complete specification of the \exit{} function which
  performs an unconditional exit of the \main{} function:
  \listinginput{1}{exit.c}
  Note that the specification of the \lstinline|may_exit|
  function is incomplete since it allows modifications of the variable
  \lstinline|status| when no exit is performed.
  Using behaviors, it is possible to distinguish between the exit case
  and the normal case, as in the following specification:
  \listinginput{8}{mayexit.c}
\end{example}

\begin{todo}
Discuter de la s�mantique du assign en pr�sence d'une terminaison abrupte
dans un contrat de fonction et de statement.
\end{todo}
% [Patrick B.] START proposition:
Contrary to \ensures clauses, \assigns\indextt{assigns}, 
\allocates\indextt{allocates} and \frees\indextt{frees} clauses of function 
and statement contracts constrain the post-state 
even when the annotated function and statement terminates respectively abruptly. 
This is shown in example~\ref{ex:assigns-and-abrupt-termination} for a function contract.
% END proposition:

\begin{todo}
  Discuter de exit versus return pour la fonction main. Discuter
    du atexit() : le post-state de la clause exit est-il avant ou
    apres les fonctions enregistrees par atexit() ?

\end{todo}

\oldremark{PC}{

  utilite de la clause assigns ?

}

\section{Dependencies information}
\label{sec:func-dep}

\experimental

An extended syntax of \assigns clauses, described in
Figure~\ref{fig:gram:dep} allows to specify data
dependencies\index{dependency} and
\emph{functional expressions}\index{functional expression}.

\begin{figure}[t]
  \fbox{\begin{minipage}{0.97\linewidth}
      \input{dependencies_modern.bnf}
    \end{minipage}}
  \caption{Grammar for dependencies information}
\label{fig:gram:dep}
\end{figure}

Such a clause indicates that the assigned values can only depend upon
the locations mentioned in the \lstinline|\from| part of the
clause. Again, this is an over-approximation: all of the locations
involved in the computation of the modified values must be present,
but some of locations might not be used in practice. If the
\lstinline|\from| clause is absent, all of the locations reachable at the
given point of the program are supposed to be used.
Moreover, for a single location, it is possible to give the precise
relation between its final value and the value of its
dependencies. This expression is evaluated in the pre-state of the
corresponding contract.

\begin{example}
  The following example is a variation over the \lstinline+array_sum+
  function in example~\ref{ex:higherorder}, in which the values of
  the array are added to a global variable \lstinline+total+.

  \listinginput{1}{sum2.c}

\end{example}


\begin{example}
 \begin{notimplementedenv}
  % This exemple is referenced in a previous section.
  The composite element modifier operators 
  can be useful for writing such functional expressions.
 \end{notimplementedenv}
  \listinginput{1}{modifier.c}
\end{example}

\section{Data invariants}
\label{sec:invariants}
\index{data invariant}\index{global invariant}\index{type invariant}
\index{invariant!data}\index{invariant!global}\index{invariant!type}
Data invariants are properties on data that are supposed to hold
permanently during the lifetime of these data. In ACSL, we distinguish
between:
\begin{itemize}
\item \emph{global} invariants and \emph{type} invariants: the former
  only apply to specified global variables, whereas the latter are
  associated to a static type, and apply to any variables of
  the corresponding type;
\item \emph{strong} invariants and \emph{weak} invariants:
  \index{invariant!strong}\index{invariant!weak}
  strong
  invariants must be valid at any time during program execution
  (more precisely at any \emph{sequence point} as defined in the C
  standard), whereas weak invariants must be valid at \emph{function
    boundaries} (function entrance and exit) but can be violated in
  between.
\end{itemize}

\begin{figure}[t]
  \fbox{\begin{minipage}{0.97\linewidth}
      \input{data_invariants_modern.bnf}
    \end{minipage}}
  \caption{Grammar for declarations of data invariants}
\label{fig:gram:datainvariants}
\end{figure}

The syntax for declaring data invariants is given in
Figure~\ref{fig:gram:datainvariants}. The strength modifier defaults
to \texttt{weak}.


\begin{example}
  In the following example, we declare
  \begin{enumerate}
  \item a weak global invariant \lstinline|a_is_positive| which specifies that
    global variable \lstinline|a| should remain positive (weakly, so
    this property might be violated temporarily between functions
    calls);
  \item a strong type invariant for variables of type \lstinline|temperature|;
  \item a weak type invariant for variables of type \lstinline|struct S|.
  \end{enumerate}
  \listinginput{1}{invariants.c}
\end{example}

\oldremark{SD8}{

  montrer un example qui exprime des proprietes d'initialisation de
  variables globales, i.e un tableau contient des valeurs distinctes 2
  a 2.

}

\subsection{Semantics}

The distinction between strong and weak invariants has to do with the
sequence points where the property is supposed to hold. The distinction
between global and type invariants has to do with the set of values on
which they are supposed to hold.

\begin{itemize}
\item Weak global invariants are properties which apply to global
  data and hold at any function entrance and function exit.

\item Strong global invariants are properties which apply to global
  data and hold at any step during execution (starting after
  initialization of these data).

\item A weak type invariant on type $\tau$ must hold at any function
  entrance and exit, and applies to any global variable or formal
  parameter with static type $\tau$. If the result of the
  function is of type $\tau$, the result must also satisfy its
  weak invariant at function exit. However, it says nothing about fields,
  array elements, memory locations, etc. of type $\tau$.

  \oldremark{Claude}{

    et les var locales d'une fonction appelante ?

  }
\item A strong type invariant on type $\tau$ must hold at any step
  during execution, and applies to any global variable, local
  variable, or formal parameter\index{formal parameter}
  with static type $\tau$. If the
  result of the function has type $\tau$, the result must also
  satisfy its strong invariant at function exit. Again, it says
  nothing about fields, array elements, memory locations, etc. of type
  $\tau$.

\end{itemize}


\begin{example}
  \notimplemented{The following example illustrates the use of a data invariant on a local static variable.}

  \listinginput{1}{out_char.c}
\end{example}



\oldremark{PC}{

  remise en cause de la politique d'invariants



Weak invariants do not hold in any state but can be temporarily violated.
Such invariants are allowed under some sound policy~\cite{barnett04jot}.

}

\oldremark{Claude}{

  il faut decider comment sont lies les declarations d'invariants et
  les type associe'.

Example:


typedef struct S {
  int x;
} S;

void reset(S *p) {
  p->x = 0;
}

//@ type invariant I(S *p) { p.x > 0 }

-> reset ne respecte pas cet invariant
}

\begin{example}
  Here is a longer example, the famous Dijkstra's Dutch flag algorithm.\\
  \listinginput{1}{flag.c}
\end{example}

\subsection{Model variables and model fields}
\index{model} \notimplemented{A \emph{model variable} is a variable
  introduced in the specification with the keyword
  \lstinline|model|}. Its type must be a logic type. Analogously,
\notimplemented {structures may have \emph{model fields}}.  These are
used to provide abstract specifications to functions whose concrete
implementation must remain private.

\begin{figure}[b]
  \fbox{\begin{minipage}{0.97\linewidth}
      \input{model_modern.bnf}
    \end{minipage}}
  \caption{Grammar for declarations of model variables and fields}
\label{fig:gram:model}
\end{figure}


The precise syntax for declaring model variables and fields is given
in Figure~\ref{fig:gram:model}. It is presented as additions to the
regular C grammar for variable declarations and structure field
declarations.

Informal semantics of model variables is as follows.
\begin{itemize}
\item Model variables can only appear in specifications. They are not
  lvalues, thus they cannot be assigned directly (unlike ghost variables,
  see below).
\item Nevertheless, a function contract might state that a model variable
  is assigned.
\item When a function contract mentions model variables:
  \begin{itemize}
  \item the precondition is implicitly existentially quantified over
    those variables;
  \item the postconditions are universally quantified over the old
    values of model variables, and existentially quantified over the new values.
  \end{itemize}
\end{itemize}
Thus, in practice, the only way to prove that a function body
satisfies a contract with model variables is to provide an invariant
relating model variables and concrete variables, as in the example
below.

\begin{example}\label{ex:model}
  Here is an example of a specification for a function which generates
  fresh integers. The contract is given in term of a model variable
  which is intended to represent the set of ``forbidden'' values,
  e.g. the values that have already been generated.
  \listinginput{1}{gen_spec_with_model.c}
  The contract is expressed abstractly, telling that
  \begin{itemize}
  \item the forbidden set of values is modified;
  \item the value returned is not in the set of forbidden values, thus
    it is ``fresh'';
  \item the new set of forbidden values contains both the value
    returned and the previous forbidden values.
  \end{itemize}
  An implementation of this function might be as follows, where a
  decision has been made to generate values in increasing order, so
  that it is sufficient to record the last value generated. This decision
  is made explicit by an invariant.
  \listinginput{1}{gen_code.c}
\end{example}

\paragraph{Remarks}

Although the syntax of model variables is close to JML model
variables, they differ in the sense that the type of a model variable
is a logic type, not a C type. Also, the semantics above is closer to
the one of B machines~\cite{abrial96:_b_book}. It has to be noticed
that program verification with model variables does not have a
well-established theoretical background~\cite{marche07,leavens07}, so
we deliberately do not provide a precise semantics in this document .

\section{Ghost variables and statements}
\label{sec:ghost}
\index{ghost}
Ghost variables and statements are like C variables and statements,
but visible only in the specifications. They are introduced by the
\lstinline|ghost| keyword at the beginning of the annotation
(i.e. \lstinline|/*@ ghost ... */| or \lstinline|//@ ghost ...|
for a one-line
ghost code, as mentioned in section~\ref{sec:gener-about-annot}).
The grammar is given in Figure~\ref{fig:gram:ghost}, in which only the
first form of annotation is used. In this figure, the \textit{C-*}
non-terminals refer to the corresponding grammar rules of the ISO standard,
without any ACSL extension. Any non terminal of the form
\textit{ghost-non-term} for which no definition is given in the figure
represents the corresponding \textit{C-non-term} entry, in which any
\textit{entry} is substituted by \textit{ghost-entry}.

The variations with respect to the C
grammar are the following:
\begin{itemize}
\item Comments must be introduced by \lstinline|//| and extend until the
  end of the line (the ghost code itself is placed inside a C
  comment. \lstinline|/* ... */| would thus lead to incorrect C code).
\item It is however possible to write multi-line annotations for ghost
  code. These annotations are enclosed between
  \notimplemented{\lstinline|/@| and \lstinline|@/|}.
  As in normal annotations, \lstinline|@|s at the beginning of a
  line and at the end of the comment (before the final \lstinline|@/|) are
  considered as blank.
  \oldremark{Claude}{
    faudrait supprimer ca. Je propose de maniere generale que le code
    annote sans les "/*@ ..*/" doit etre parsable aussi, meme s'il
    n'est pas parsable par un compilo.
}
\item \notimplemented{Logical types, such as \lstinline|integer| or
\lstinline|real| are authorized in ghost code}.
\item \notimplemented{A non-ghost function can take ghost
    parameters}.
  If such a ghost
  clause is present in the declarator, then the list of ghost
  parameters must be non-empty and fixed (no vararg ghost). The call
  to the function must then provide the appropriate number of ghost parameters.
\item Any non-ghost \textit{if-statement} which does not have a non-ghost
  \lstinline|else| clause \notimplemented{can be augmented with a ghost
    one}. Similarly, a non-ghost
  \lstinline|switch| can have \notimplemented{a ghost \lstinline|default:|
    clause}
  if it does not have a non-ghost one
  (there are however semantical restrictions for valid
  ghost labelled statements in a switch, see next paragraph for details).
\end{itemize}

\begin{figure}[t]
  \fbox{\begin{minipage}{0.98\linewidth}
      \input{ghost_modern.bnf}
    \end{minipage}}
  \caption{Grammar for ghost statements}
\label{fig:gram:ghost}
\end{figure}


\paragraph{Semantics of Ghost Code}
\label{sec:semantics-ghost-code}
The question of semantics is essential for ghost code.
\notimplemented[Not checked in the current implementation]{Informally, the
semantics requires that ghost statements do not change the regular
program execution} This implies several conditions, including e.g:
\begin{itemize}
\item Ghost code cannot modify a non-ghost C variable.
\item Ghost code cannot modify a non-ghost structure field.
\item If \lstinline|p| is a ghost pointer pointing to a non-ghost
  memory location, then it is forbidden to assign \lstinline|*p|.
\item Body of a ghost function is ghost code, hence do not modify
  non-ghost variables or fields.
\item If a non-ghost C function is called in ghost code, it must not
  modify non-ghost variables or fields.
\item If a structure has ghost fields, the \lstinline|sizeof| of the
  structure is the same has the structure without ghost fields. Also,
  alignment of fields remains unchanged.
\item The control-flow graph of a function must not be altered by
  ghost statements. In particular, no ghost \lstinline|return| can appear
  in the body of a non-ghost function. Similarly, ghost
  \lstinline|goto|, \lstinline|break|, and continue
  \lstinline|continue| cannot jump
  outside of the innermost non-ghost enclosing block.
\end{itemize}

Semantics is specified as follows. First, the
execution of a program with ghost code involves a \emph{ghost memory heap}
and a \emph{ghost stack}, disjoint from the regular heap and stack.
Ghost variables lie in the ghost heap, so as the ghost field of
structures. Thus, every memory side-effect can be classified as ghost
or non-ghost. Then, the semantics is that memory side-effects of ghost
code must always be in the ghost heap or the ghost stack.

Notice that this semantics is not statically decidable.
\notimplemented{It is left to tools to provide approximations},
correct in the sense that any code
statically detected as ghost must be semantically ghost.

\begin{example}
  The following example shows some invalid assignments of ghost pointers:
  \listinginput{1}{ghostpointer.c}
\end{example}

\begin{example}
  The following example shows some invalid ghost statements:
  \listinginput{1}{ghostcfg.c}
\end{example}

\oldremark{Claude}{
  d'autres exemples, valides ou non valides, serait bienvenus
}

\oldremark{All}{

On veut des tableaux ghost de taille infinie -> cf tableaux fonctionnels de la section sur les assigns from

}

\paragraph{Differences between model variables and ghost variables}


A ghost variable is an additional specification variable which is
assigned in ghost code like any C variable. On the other hand, a model
variable cannot be assigned, but one can state it is modified and can
express properties about the new value, in a non-deterministic way,
using logic assertions and invariants.
In other words, specifications using ghost variable assignments
are executable.

\begin{example}
  \label{ex:gen_code}
  The example~\ref{ex:model} can also be specified with a ghost
  variable instead of a model variable:

  \listinginput{1}{gen_spec_with_ghost.c}

\end{example}

\subsection{Volatile variables}\label{sec:volatile-variables}
\index{volatile}
\experimental

Volatile variables can not be used in logic terms, since reading such
a variable may have a side effect, in particular two successive reads
may return different values.

\begin{figure}[h]
  \fbox{\begin{minipage}{0.97\linewidth}
      \input{volatile-gram_modern.bnf}
    \end{minipage}}
  \caption{Grammar for volatile constructs}
\label{fig:gram:volatile}
\end{figure}

Specifying properties of a volatile variable may be done via a
specific construct to attach two ghost functions to it. This
construct, described by the grammar of Figure~\ref{fig:gram:volatile}, has
the following shape:
\begin{lstlisting}[style=c-basic,firstnumber=1,name=volatile,numbers=left]
volatile $\tau$ x;
\end{lstlisting}
\begin{notimplementedenv}
\begin{lstlisting}[style=c-basic,name=volatile,numbers=left]
//@ volatile x reads f writes g;
\end{lstlisting}
\end{notimplementedenv}
where \lstinline|f| and \lstinline|g| are ghost functions with the following prototypes:
\begin{lstlisting}[style=c-basic,name=volatile,numbers=left]
$\tau$ f(volatile $\tau$* p);
$\tau$ g(volatile $\tau$* p, $\tau$ v);
\end{lstlisting}
This must be understood as a special construct to instrument the C
code, where each access to the variable \lstinline|x| is replaced by a call
to \lstinline|f(&x)|, and each assignment to \lstinline|x| of a value
\lstinline|v|
is replaced by \lstinline|g(&x,v)|. If a given volatile variable is only read 
or only written to, the unused accessor function can be omitted from the 
\lstinline|volatile| construct.

\begin{example}
  The following code is instrumented in order to inject fixed values
  at each read of variable \lstinline|x|, and collect written values.
  \listinginput{1}{volatile.c}
\end{example}

\section{Undefined values, dangling pointers}

\subsection{Initialization}
\indexttbs{initialized}
\lstinline|\initialized| is a predicate taking a set of l-values as
argument and means that each l-value in this set is initialized.

\begin{example}
  In the following, the assertion is true.
  \listinginput{1}{initialized.c}
\end{example}

\subsection{Unspecified values}
\indexttbs{specified}
\begin{notimplementedenv}\lstinline|\specified|\end{notimplementedenv}
is a predicate taking a set of l-values as
argument and means that each l-value in this set has a \emph{specified
  value}: its value is not a dangling pointer (that is, the value
is not the address of a local variable refered to outside of its scope)

\begin{example}
  In the following, the assertion is not true.
  \listinginput{1}{specified.c}
\end{example}



%\section{Module constructions}
%
%\experimental
%
%how to encapsulate several functions...


%\section{Attribute annotations}
%\label{sec:attribute_annot}
%
%\experimental
%
%These are annotations allowing to add attributes on variables, like regular C attributes (const, volatile, restrict).

% The purpose of these are various:
% \begin{itemize}
% \item specific verifications, like \url{http://www.cs.umd.edu/~jfoster/cqual}
% \item specific analyses, like region
%   analyses~\cite{hubert07hav} a la Cyclone~\cite{grossman02pldi}.
% \item allowing to attach model variables/fields to types that are not
%   structures~\cite{filliatr07queens}.
% \end{itemize}

%%% Local Variables:
%%% mode: latex
%%% TeX-PDF-mode: t
%%% TeX-master: "main"
%%% ispell-local-dictionary: "english"
%%% End:


\chapter{Libraries}
\label{chap:lib}

\section{Libraries of logic specifications}
\label{sec:speclibraries}

\subsection{Real numbers}
\label{sec:libreal}

\subsection{Finite lists}


\subsection{Sets and Maps}

\section{Jessie library: logical adressing of memory blocks}
\label{sec:jessie}

\subsection{Abstract level of pointer validity}

\subsection{Strings}

\experimental

\section{Memory leaks}

\experimental



\chapter{Conclusion}

This document presents a Behavioral Interface Specification Language
for ANSI C source code. It provides a common basis that can be
shared among different tools.
The specification language described here is intended to evolve in the
future and remain open to additional constructions.
One interesting possible extension regards ``temporal''
properties in a large sense, such as liveness properties, which can
sometimes be simulated by regular specifications with ghost
variables~\cite{giorgetti06fase}, or properties on evolution of data
over the time, such as the history constraints of JML, or in the Lustre
assertion language.

%%% Local Variables:
%%% mode: latex
%%% TeX-PDF-mode: t
%%% TeX-master: "main"
%%% End:


\appendix

\chapter{Appendices}
\label{chap:appendix}

%\section{Quick reference card}

%\todo{by Benjamin}

\section{Glossary}
\label{sec:glossary}

\begin{description}
\item[pure expressions] \index{pure expression} In ACSL setting, a
  \emph{pure} expression is a C expression which contains no assignments, no
  incrementation operator \lstinline|++| or \lstinline|--|, no function call,
  and no access to a volatile object. The set of pure expression is a
  subset of the set of C expressions without side effect (C
  standard~\cite{KR88,standardc99}, \S 5.1.2.3, alinea 2).

\item[left-values] \index{left-value}\index{lvalue}

  A \emph{left-value} (\emph{lvalue} for short) is an expression which
  denotes some place in the memory during program execution, either on
  the stack, on the heap, or in the static data segment. It can be
  either a variable identifier or an expression of the form \lstinline|*e|,
  \lstinline|e[e]|, \lstinline|e.id| or \lstinline|e->id|, where
  \lstinline |e| is any expression and \lstinline|id| a field name.
  See C standard, \S 6.3.2.1 for a more
  detailed description of lvalues.

  A \emph{modifiable lvalue} is an lvalue allowed in the left part of
  an assignment. In essence, all lvalues are modifiable except
  variables declared as \texttt{const} or of some array type with
  explicit length.

  \oldremark{DP}{
De'tailler (avec un exemple), meme si c'est dans la
    norme, la difference entre lvalue modifiable et non-modifiable
  }


\item[pre-state and post-state]
    \index{pre-state}\index{post-state}

    For a given function call, the \emph{pre-state} denotes the
    program state at the beginning of the call, including the
    current values for the function parameters. The \emph{post-state}
    denotes the program state at the return of the call.

\item[function behavior] \index{function behavior} \index{behavior}

  A \emph{function behavior} (\emph{behavior} for short) is a set of
  properties relating the pre-state and the post-state for a
  possibly restricted set of pre-states (behavior \emph{assumptions}).

\item[function contract] \index{function contract} \index{contract}

  A \emph{function contract} (\emph{contract} for short) forms a
  specification of a function, consisting of the combination of a
  precondition (a requirement on the pre-state for any caller to that
  function), a collection of behaviors, and possibly a measure in case
  of a recursive function.

\end{description}

\section{Comparison with JML}
\label{sec:comp-jml}

Although we took our inspiration in the Java Modeling Language (aka
JML~\cite{jmlhomepage}), ACSL is notably different from JML in
two crucial aspects:

\begin{itemize}
\item ACSL is a BISL for C, a low-level structured language, while JML
  is a BISL for Java, an object-oriented inheritance-based high-level
  language. Not only the language features are not the same but the
  programming styles and idioms are very different, which entails also
  different ways of specifying behaviors. In particular, C has no
  inheritance nor exceptions, and no language support for the simplest
  properties on memory (\emph{e.g}, the size of an allocated memory block).
\item JML relies on runtime assertion checking (RAC) when typing,
  static analysis and automatic deductive verification fail. The
  example of CCured~\cite{necula02ccured,condit03ccured}, that adds
  strong typing to C by relying on RAC too, shows that it is not possible
  to do it in a modular way. Indeed, it is necessary to modify the
  layout of C data structures for RAC, which is not modular. The
  follow-up project Deputy~\cite{condit07deputy} thus reduces the
  checking power of annotations in order to preserve modularity.  On
  the contrary, we choose not to restrain the power of annotations
  (\emph{e.g.}, all first order logic formulas are allowed). To that end, we
  rely on manual deductive verification using an interactive theorem
  prover (\emph{e.g.}, Coq) when every other technique failed.
\end{itemize}

\noindent
In the remainder of this chapter, we describe these differences in
further details.

\subsection{Low-level language vs. inheritance-based one}

\subsubsection*{No inherited specifications}

JML has a core notion of inheritance of specifications, that duplicates
in specifications the inheritance feature of Java.  Inheritance
combined with visibility and modularity account for a number of
complex features in JML (\emph{e.g}, \verb|spec_public| modifier, data
groups, represents clauses, etc), that are necessary to express the
desired inheritance-related specifications while respecting visibility
and modularity. Since C has no inheritance, these intricacies are
avoided in ACSL.

\subsubsection*{Error handling without exceptions}
\label{sec:errorhandling}

The usual way of signaling errors in Java is through
exceptions. Therefore, JML specifications are tailored to express
exceptional postconditions, depending on the exception raised. Since C
has no exceptions, ACSL does not use exceptional
specifications. Instead, C programmers are used to signal errors by
returning special values, like mandated in various ways in the C standard.

\begin{example}
In \S 7.12.1 of the standard, it is said that functions in <math.h>
signal errors as follows:
``On a domain error, [...] the integer expression errno
acquires the value EDOM.''
\end{example}

\begin{example}
In \S 7.19.5.1 of the standard, it is said that function fclose signals
errors as follows:
``The fclose function returns [...] EOF if any errors were detected.''
\end{example}

\begin{example}
In \S 7.19.6.1 of the standard, it is said that function fprintf
signals errors as follows:
``The fprintf function returns [...] a negative value if an output or
encoding error occured.''
\end{example}

\begin{example}
In \S 7.20.3 of the standard, it is said that memory management functions
signal errors as follows:
``If the space cannot be allocated, a null pointer is returned.''
\end{example}

As shown by these few examples, there is no unique way to signal
errors in the C standard library, not mentioning user-defined
functions. But since errors are signaled by returning special values, it
is sufficient to write an appropriate postcondition:

\begin{listing-nonumber}
/*@ ensures \result == error_value || normal_postcondition; */
\end{listing-nonumber}

% \noindent
% A tool could easily set error conditions aside, by providing an
% appropriate extension of behaviors, \emph{e.g}, using a new keyword
% \verb|failswith|:

% \begin{flushleft}\ttfamily
% /*@ failswith $\mathit{\result == error\_value}$; \\
% ~~@ ensures $\mathit{normal\_postcondition}$; \\
% ~~@*/
% \end{flushleft}


%\input{fwrite-malloc.pp}


\subsubsection*{C contracts are not Java ones}

In Java, the precondition of the following function that nullifies an
array of characters is always true. Even if there was a precondition
on the length of array {\ttfamily a}, it could easily be expressed using
the Java expression {\ttfamily a.length} that gives the dynamic length
of array {\ttfamily a}.

\begin{listing}{1}
public static void Java_nullify(char[] a) {
  if (a == null) return;
  for (int i = 0; i < a.length; ++i) {
    a[i] = 0;
  }
}
\end{listing}

On the contrary, the precondition of the same function in C, whose
definition follows, is more involved. First, remark that the
C programmer has to add an extra argument for the size of the array,
or rather a lower bound on this array size.

\begin{listing}{1}
void C_nullify(char* a, unsigned int n) {
  int i;
  if (n == 0) return;
  for (i = 0; i < n; ++i) {
    a[i] = 0;
  }
}
\end{listing}

\noindent
A correct precondition for this function is the following:

\begin{listing-nonumber}
/*@ requires \valid(a + 0..(n-1)); */
\end{listing-nonumber}

where predicate \valid is the one defined in Section~\ref{subsec:memory}.
(note that \lstinline|\valid(a + 0..(-1))| is the same as
\lstinline|\valid(\empty)| and thus is true regardless of the validity of
\lstinline|a| itself).
When \lstinline|n| is null, \lstinline|a| does
not need to be valid at all, and when \lstinline|n| is strictly
positive, \lstinline|a| must point to an array of size at least
\lstinline|n|. To make it more obvious, the C programmer adopted a
defensive programming style, which returns immediately when \lstinline|n| is
null. We can duplicate this in the specification:

\begin{listing-nonumber}
/*@ requires n == 0 || \valid(a + 0..(n-1)); */
\end{listing-nonumber}

Usually, many memory requirements are only necessary for some paths
through the function, which correspond to some particular
behaviors, selected according to some tests performed along the
corresponding paths. Since C has no memory
primitives, these tests involve other variables that the C programmer
added to track additional information, like {\ttfamily n} in our example.

To make it easier, it is possible in ACSL to distinguish between the
\lstinline|assumes| part of a behavior, that specifies the tests that need
to succeed for this behavior to apply, and the \lstinline|requires| part
that specifies the additional preconditions that must be true when a
behavior applies. The specification for our example can then be
translated into:

\begin{listing}{1}
/*@ behavior n_is_null:
  @   assumes n == 0;
  @ behavior n_is_not_null:
  @   assumes n > 0;
  @   requires \valid(a + 0..(n-1));
  @*/
\end{listing}

This is equivalent to the previous requirement, except here behaviors
can be completed with postconditions that belong to one behavior only.
Contrary to JML, the set of behaviors for a function do not
necessarily cover all cases of use for this function, as mentioned in
Section~\ref{subsec:behaviors}. This allows for partial
specifications, whereas JML behaviors cannot offer such
flexibility. Here, Our two behaviors are clearly mutually exclusive,
and, since \lstinline|n| is an \lstinline|unsigned int|, our
they cover all the possible cases. We could have specified that as well, by
adding the following lines in the contract (see
Section~\ref{sec:compl-behav}).
\begin{listing}{last}
  @ ...
  @ disjoint behaviors;
  @ complete behaviors;
  @*/
\end{listing}

\subsubsection*{ACSL contracts vs. JML ones}

To fully understand the difference between specifications in ACSL and
JML, we detail in below the requirement on the pre-state and
the guarantee on the post-state given by behaviors in JML and ACSL.

A JML contract is either \emph{lightweight} or \emph{heavyweight}.
For the purpose of our comparison, it is sufficient to know that a
lightweight contract has \lstinline|requires| and \lstinline|ensures| clauses
all at the same level, while an heavyweight contract has multiple
behaviors, each consisting of \lstinline|requires| and \lstinline|ensures|
clauses. Although it is not possible in JML to mix both styles, we can
define here what it would mean to have both, by conjoining the
conditions on the pre- and the post-state.
Here is an hypothetical JML contract mixing lightweight and
heavyweight styles:
\begin{listing}{1}
/*@ requires $P_1$;
  @ requires $P_2$;
  @ ensures  $Q_1$;
  @ ensures  $Q_2$;
  @ behavior $x_1$:
  @   requires $A_1$;
  @   requires $R_1$;
  @   ensures $E_1$;
  @ behavior $x_2$:
  @   requires $A_2$;
  @   requires $R_2$;
  @   ensures $E_2$;
  @*/
\end{listing}
It assumes from the pre-state the condition:
\begin{listing-nonumber}
$P_1$ && $P_2$ && (($A_1$ && $R_1$) || ($A_2$ && $R_2$))
\end{listing-nonumber}
and guarantees that the following condition holds in post-state:
\begin{listing-nonumber}
$Q_1$ && $Q_2$ &&
  (\old($A_1$ && $R_1$) ==> $E_1$) && (\old($A_2$ && $R_2$) ==> $E_2$)
\end{listing-nonumber}


Here is now an ACSL specification:

\begin{listing}{1}
/*@ requires $P_1$;
  @ requires $P_2$;
  @ ensures  $Q_1$;
  @ ensures  $Q_2$;
  @ behavior $x_1$:
  @   assumes $A_1$;
  @   requires $R_1$;
  @   ensures $E_1$;
  @ behavior $x_2$:
  @   assumes $A_2$;
  @   requires $R_2$;
  @   ensures $E_2$;
  @*/
\end{listing}

\noindent
Syntactically, the only difference with the JML specification is the
addition of the \lstinline|assumes| clauses.
Its translation to assume-guarantee is however quite different.
It assumes from the pre-state the condition:

\begin{listing-nonumber}
  $P_1$ && $P_2$ && ($A_1$ ==> $R_1$) && ($A_2$ ==> $R_2$)
\end{listing-nonumber}

\noindent
and guarantees that the following condition holds in the post-state:

\begin{listing-nonumber}
$Q_1$ && $Q_2$ && (\old($A_1$) ==> $E_1$) && (\old($A_2$) ==> $E_2$)
\end{listing-nonumber}

Thus, ACSL allows to distinguish between the clauses that control
which behavior is active (the \lstinline|assumes| clauses) and the
clauses that are preconditions for a particular behavior (the internal
\lstinline|requires| clauses). In addition, as mentioned above, there is
by default no requirement in ACSL for the specification to be complete (The
last part of the JML condition on the pre-state). If desired, this has
to be precised explicitly with a \lstinline+complete behaviors+ clause as
seen in Section~\ref{sec:compl-behav}.

\subsection{Deductive verification vs. RAC}

\subsubsection*{Sugar-free behaviors}

As explained in details in~\cite{raghavan00desugaring}, JML
heavyweight behaviors can be viewed as syntactic sugar (however
complex it is) that can be translated automatically into more basic
contracts consisting mostly of pre- and postconditions and frame
conditions.  This allows complex nesting of behaviors from the user
point of view, while tools only have to deal with basic contracts. In
particular, the major tools on JML use this desugaring process, like
the Common JML tools to do assertion checking, unit testing,
etc. (see~\cite{leavens00jml}) and the tool ESC/Java2 for
automatic deductive verification of JML specifications
(see~\cite{kiriny07implnotes}).

One issue with such a desugaring approach is the complexity of the
transformations involved, as \emph{e.g.} for desugaring assignable clauses
between multiple \textit{spec-cases} in
JML~\cite{raghavan00desugaring}.  Another issue is precisely that
tools only see one global contract, instead of multiple independent
behaviors, that could be analyzed separately in more detail.
Instead, we favor the view that a function implements multiple
behaviors, that can be analyzed separately if a tool feels like
it. Therefore, we do not intend to provide a desugaring process.

\subsubsection*{Axiomatized functions in specifications}

JML only allows pure Java methods to be called in
specifications~\cite{leavens00preliminary}. This is certainly
understandable when relying on RAC: methods called should be defined
so that the runtime can call them, and they should not have
side-effects in order not to pollute the program they are supposed to
annotate.

In our setting, it is desirable to allow calls to logical functions in
specifications. These functions may be defined, like program ones, but
they may also be only declared (with a suitable declaration of \reads
clause) and defined through an axiomatization.
This makes for richer specifications that may be useful either in
automatic or in manual deductive verification.

\subsection{Syntactic differences}

The following table summarizes the difference between JML keywords and
ACSL ones, when the intent is the same, although minor differences
might exist.
\begin{center}
\begin{tabular}{|l|l|}
\hline
  JML                  & ACSL \\ \hline
  modifiable,assignable           & assigns \\
  measured\_by         & decreases \\
  loop\_invariant      & loop invariant \\
  decreases            & loop variant \\
  \lstinline|(\forall $\tau$ x ; P ; Q)| &
       \lstinline|(\forall $\tau$ x ; P ==> Q)| \\
  \lstinline|(\exists $\tau$ x ; P ; Q)| &
        \lstinline|(\exists $\tau$ x ; P && Q)| \\
  \lstinline|\max $\tau$ x ; a <= x <= b ; f)| &
        \lstinline|\max(a,b,\lambda $\tau$ x ; f)| \\
  \hline
\end{tabular}
\end{center}



\oldremark{PC}{

  ne pas diverger de JML quand ce n'est pas justifie'


  -> rediscuter du decreases versus loop variant

}


%%% Local Variables:
%%% mode: latex
%%% TeX-PDF-mode: t
%%% TeX-master: "main"
%%% End:


\section{Typing rules}
\label{sec:typingrules}

Disclaimer: this section is unfinished, it is left here just to give an idea of what it will look like when completed.

\subsection{Rules for terms}

Integer promotion:
\[
\frac{\Gamma \vdash e : \tau}{\Gamma \vdash e : \integer}
\]
if $\tau$ is any C integer type \verb|char|, \verb|short|, \verb|int|, or \verb|long|, whatever attribute they have, in particular signed or unsigned

Variables:
\[
\frac{}{\Gamma \vdash id : \tau} \mbox{ if $id:\tau\in\Gamma$}
\]

Unary integer operations:
\[
\frac{\Gamma \vdash t : \integer}{\Gamma \vdash op~t : \integer} \mbox{ if $op\in \{+,-,\sim\}$}
\]

Boolean negation:
\[
\frac{\Gamma \vdash t : \bool}{\Gamma \vdash !~t : \bool}
\]

Pointer dereferencing:
\[
\frac{\Gamma \vdash t : \tau*}{\Gamma \vdash *t : \tau}
\]

Address operator:
\[
\frac{\Gamma \vdash t : \tau}{\Gamma \vdash \&t : \tau*}
\]

Binary
\[
\frac{\Gamma \vdash t_1 : \integer\qquad\Gamma \vdash t_2 : \integer}{\Gamma \vdash t_1~op~t_2 : \integer} \mbox{ if $op\in \{+,-,*,/,\%\}$}
\]
\[
\frac{\Gamma \vdash t_1 : \real\qquad\Gamma \vdash t_2 : \real}{\Gamma \vdash t_1~op~t_2 : \real} \mbox{ if $op\in \{+,-,*,/\}$}
\]
\[
\frac{\Gamma \vdash t_1 : \integer\qquad\Gamma \vdash t_2 : \integer}{\Gamma \vdash t_1~op~t_2 : \bool} \mbox{ if $op\in \{==,!=,<=,<,>=,>\}$}
\]
\[
\frac{\Gamma \vdash t_1 : \real\qquad\Gamma \vdash t_2 : \real}{\Gamma \vdash t_1~op~t_2 : \bool} \mbox{ if $op\in \{==,!=,<=,<,>=,>\}$}
\]
\[
\frac{\Gamma \vdash t_1 : \tau*\qquad\Gamma \vdash t_2 : \tau*}{\Gamma \vdash t_1~op~t_2 : \bool} \mbox{ if $op\in \{==,!=,<=,<,>=,>\}$}
\]

(to be continued)


\subsection{Typing rules for sets}
%VP: A-t-on besoin d'un logic environment?
We consider the typing judgement $\Gamma,\Lambda \vdash s : \tau,b$
meaning that $s$ is a set of terms of type $\tau$, which is moreover a
set of locations if the boolean $b$ is true.
$\Gamma$ is the C environment and $\Lambda$ is the logic environment.

Rules:
\[
\frac{}{\Gamma,\Lambda \vdash id : \tau,true} \mbox{ if $id:\tau \in \Gamma$}
\]
\[
\frac{}{\Gamma,\Lambda \vdash id : \tau,true} \mbox{ if $id:\tau \in \Lambda$}
\]
\[
\frac{\Gamma,\Lambda\vdash s:\tau*,b}{\Gamma,\Lambda \vdash *s: \tau,true}
\]
%VP: cette regle ne veut rien dire.
%CM maintenant si, si l'environnement donne les types des champs
\[
\frac{id:\tau \quad s: set<struct~S*>}{\vdash s->id : set<\tau>}
\]
\[
\frac{\Gamma,b\cup \Lambda \vdash e: tset \tau}
{\Gamma,\Lambda\vdash \{ e \mid b ; P \} : tset \tau }
\]
\[
\frac{\Gamma,\Lambda\vdash e_1:\tau,b \quad \Gamma,\Lambda\vdash e_2:\tau,b}
{\Gamma,\Lambda\vdash e_1,e_2: \tau,b}
\]

%%% Local Variables:
%%% mode: latex
%%% TeX-PDF-mode: t
%%% TeX-master: "main"
%%% End:


\section{Specification Templates}\label{sec:spec-templ}
This section describes some common issues that may occur when writing
an ACSL specification and proposes some solution to overcome them

\subsection{Accessing a C variable that is masked}

The situation may happen where it is necessary to refer
in an annotation to a C variable that is masked at that point.
For instance,
a function contract may need to refer to a global variable that has the
same name as a function parameter, as in the following code:

\listinginput{1}{glob_var_masked.c}

In order to write the \verb|assigns| clause for \verb|f|, we must
access the global variable \verb|x|, since \verb|f| calls \verb|g|,
which can modify \verb|x|. This is not possible with C scoping rules,
as \verb|x| refers to the parameter of \verb|f| in the scope of the
function.

A solution is to use a ghost pointer to \verb|x|, as shown in
the following code:

\listinginput{1}{glob_var_masked_sol.c}


%%% Local Variables:
%%% mode: latex
%%% TeX-PDF-mode: t
%%% TeX-master: "main"
%%% End:


\section{Illustrative example}

This is an attempt to define an example for ACSL, much as the Purse
example in JML description papers.  It is a memory allocator, whose
main functions are \lstinline|memory_alloc| and \lstinline|memory_free|, to
respectively allocate and deallocate memory.  The goal is to exercise
as much as possible of ACSL.

%\input{acsl_allocator.pp}
\listinginput{1}{acsl_allocator.c}

\section{Changes}

\subsection{Version \version}

\subsection{Version 1.5}

\begin{itemize}
\item Clarify the status of \lstinline|loop invariant| in presence of
  \lstinline|break| or side-effects in the loop test.
\item Introduction of \lstinline|\with| keyword for functional updates.
\item Added bnf entry for completeness of function behaviors.
\item Order of clauses in statement contracts is now fixed.
\item Requires clauses are allowed before behaviors of statement contracts.
\item Added explicit singleton construct for sets.
\item Introduction of logical arrays.
\item Operations over pointers and arrays have been precised.
\end{itemize}

\subsection{Version 1.4}
\begin{itemize}
\item Added UTF-8 counterparts for built-in types
  (\lstinline|integer|, \lstinline|real|, \lstinline|boolean|).
\item Fixed typos in the examples corresponding to features
  implemented in Frama-C.
\item Order of clauses in function contracts is now fixed.
\item Introduction of abrupt termination clauses.
\item Introduction of \lstinline|axiomatic| to gather predicates, logic
  functions, and their defining axioms.
\item Added specification templates appendix for common specification issues.
\item Use of sets as first-class term has been precised.
\item Fixed semantics of predicate \lstinline|\separated|.
\end{itemize}

\subsection{Version 1.3}
\begin{itemize}
\item Functional update of structures.
\item Terminates clause in function behaviors.
\item Typos reported by David Mentr\'e.
\end{itemize}

\subsection{Version 1.2}
This is the first public release of this document.

\cleardoublepage
\addcontentsline{toc}{chapter}{\bibname}
\bibliographystyle{plain}
%\IfFileExists{biblio-demons.tex}{%
%\input{biblio-demons.tex}
%}{
\bibliography{./biblio}
%}

\cleardoublepage
\addcontentsline{toc}{chapter}{\listfigurename}
\listoffigures

\cleardoublepage
\addcontentsline{toc}{chapter}{\indexname}
\printindex


\end{document}

%%% Local Variables:
%%% mode: latex
%%% TeX-PDF-mode: t
%%% TeX-master: t
%%% ispell-local-dictionary: "english"
%%% End:
